\documentclass[11pt,reqno]{amsart}

%% ---------------------------------------------------------------------------
%%  No balanced 5-colouring of K_26: the Erdos-Gyarfas conjecture for r = 5
%%  Erdos Problem #617, first open case.
%%  Compile: pdflatex erdos617-r5.tex  (twice, for cross-references)
%% ---------------------------------------------------------------------------

\usepackage[T1]{fontenc}
\usepackage[utf8]{inputenc}
\usepackage{amsmath,amssymb,amsthm}
\usepackage{mathtools}
\usepackage[margin=1.1in]{geometry}
\usepackage{microtype}
\usepackage{enumitem}
\usepackage{booktabs}
\usepackage{array}
\usepackage[colorlinks=true,linkcolor=blue!55!black,citecolor=green!45!black,urlcolor=blue!60!black]{hyperref}
\usepackage{xcolor}

%% ---- theorem environments -------------------------------------------------
\newtheorem{theorem}{Theorem}[section]
\newtheorem{lemma}[theorem]{Lemma}
\newtheorem{proposition}[theorem]{Proposition}
\newtheorem{corollary}[theorem]{Corollary}
\newtheorem{conjecture}[theorem]{Conjecture}
\theoremstyle{definition}
\newtheorem{definition}[theorem]{Definition}
\newtheorem{remark}[theorem]{Remark}
\newtheorem{fact}[theorem]{Fact}
\theoremstyle{remark}
\newtheorem*{claim}{Claim}

%% ---- named theorems for the introduction ----------------------------------
\newtheorem{maintheorem}{Theorem}
\renewcommand{\themaintheorem}{\Alph{maintheorem}}

%% ---- macros ---------------------------------------------------------------
\newcommand{\N}{\mathbb{N}}
\newcommand{\F}{\mathbb{F}}
\newcommand{\eq}[1]{\eqref{#1}}
\newcommand{\abs}[1]{\left|#1\right|}
\newcommand{\ee}{\mathrm{e}}
\newcommand{\ex}{\operatorname{ex}}
\newcommand{\capeleven}{\textsf{cap-11}}
\DeclareMathOperator{\Ncap}{N}
\DeclareMathOperator{\ed}{e}

\begin{document}

\title[Erd\H{o}s Problem 617 at $r=5$: a candidate resolution]{A candidate resolution of Erd\H{o}s Problem 617 for $r=5$:\\
no balanced 5-colouring of $K_{26}$, with machine verification}

\author{}
\date{\today}

\thanks{This document is a verification-ready report, not an authored announcement.
The mathematics was produced and internally reviewed by AI systems; it is offered
here for independent verification by mathematicians. See
Section~\ref{sec:howto} for a verification guide and Section~\ref{sec:provenance}
for a full account of provenance and what has and has not been checked.
\emph{Version~1.1 (2026-07-12):} the Lean development now carries no mathematical
hypothesis --- the Kang--Pikhurko equality classification at $(5,21)$, on which
version~1.0 (earlier the same day) was conditional, has been discharged in Lean
(Section~\ref{ssec:lean-brouwer}), at the cost of ten added \texttt{native\_decide}
construction axioms (Section~\ref{sec:lean}).}

\begin{abstract}
A $5$-colouring of the edges of the complete graph $K_{26}$ is \emph{balanced} if
every $6$ of its vertices span all $5$ colours. Erd\H{o}s and Gy\'arf\'as
conjectured in 1999 that for every $r\ge 3$ the complete graph $K_{r^2+1}$ admits
no balanced $r$-colouring; this is Erd\H{o}s Problem~$617$. They proved it for
$r=3$ and $r=4$, noted that it fails for $r=2$, and observed that the
corresponding statement for $K_{r^2}$ is false for infinitely many $r$, so any
proof must use the ``$+1$'' essentially. We present a \emph{candidate resolution}
of the first open case: a proof that there is no balanced $5$-colouring of
$K_{26}$. Equivalently, the largest complete graph with a balanced $5$-colouring is
$K_{25}$, realised by the affine plane $\mathrm{AG}(2,5)$.

The proof reduces, by an elementary vertex-deletion argument, to two statements
about a single colour class on $25$ vertices, and these are established by a
self-contained double-counting recursion driven by the ``every $6$ vertices span
at most $11$ edges of one colour'' constraint, one classical extremal input
(Brouwer's non-partite Tur\'an bound with the Kang--Pikhurko equality
classification), and four finite graph facts certified by SAT and independently
DRAT-checked. As a by-product we prove that the merged affine-plane colouring of
$K_{q^2}$ admits no balanced one-vertex extension, for every prime power
$q\ge 3$, so any counterexample to the conjecture must be non-affine.

The entire argument has been formalised in Lean~4, and the headline theorem
\texttt{erdos\_617\_r5\_\allowbreak unconditional} carries \emph{no} mathematical hypothesis: it is
\texttt{sorry}-free and unconditional in the exact sense that it rests only on a
disclosed trust base --- the three standard Lean axioms together with
\texttt{native\_decide} compiler-level reflection for four SAT certificates and a few
finite graph-construction facts. The Kang--Pikhurko equality classification at $(5,21)$, on
which the result was previously conditional, is now itself proved in Lean (as is
Brouwer's edge bound), though a conditional form assuming that classification is
retained for separate audit. This
document is a verification-ready report rather than an authored announcement: the mathematics was produced and internally reviewed by AI systems and
is offered here for independent verification. Section~\ref{sec:howto} gives a
verification guide (a $30$-minute machine-checkable path and a $120$-minute reading
path); Section~\ref{sec:provenance} discloses the role of the AI systems in full and
states plainly what has and has not been checked.
\end{abstract}

\maketitle

\setcounter{tocdepth}{1}
%\tableofcontents

%% ===========================================================================
\section{Introduction}
%% ===========================================================================

\subsection{The conjecture}
For an integer $r\ge 2$, call an edge-colouring
$\chi\colon E(K_n)\to\{0,1,\dots,r-1\}$ of the complete graph $K_n$
(not necessarily using all colours) \emph{balanced} if every set of $r+1$
vertices sees all $r$ colours among the $\binom{r+1}{2}$ edges it spans. Writing
$G_c$ for the graph on $V(K_n)$ whose edges are the colour-$c$ edges (the
\emph{colour class} $c$), balancedness is equivalent to the statement that every
$G_c$ has independence number $\alpha(G_c)\le r$: an $(r+1)$-set missing colour
$c$ is exactly an independent set of size $r+1$ in $G_c$.

Erd\H{o}s and Gy\'arf\'as introduced balanced colourings in their study of split
colourings~\cite{ErGy99} and made the following conjecture (their Conjecture~1,
the $n=2$ case of a more general question).

\begin{conjecture}[Erd\H{o}s--Gy\'arf\'as~\cite{ErGy99}; Erd\H{o}s Problem
$617$~\cite{Bloom617}]\label{conj:main}
For every $r\ge 3$, the complete graph $K_{r^2+1}$ has no balanced $r$-colouring.
Equivalently, if the edges of $K_{r^2+1}$ are $r$-coloured, then some $r+1$
vertices span an induced $K_{r+1}$ in which at least one colour is missing.
\end{conjecture}

It is convenient to set
\[
  \Ncap(r)\;=\;\max\{\, n : K_n \text{ has a balanced } r\text{-colouring}\,\}.
\]
Since deleting a vertex cannot increase any independence number, a balanced
colouring of $K_n$ restricts to a balanced colouring of $K_m$ for every
$r+1\le m\le n$; thus $\Ncap(r)$ is well defined and Conjecture~\ref{conj:main}
is precisely the assertion $\Ncap(r)\le r^2$ for all $r\ge 3$.

\subsection{What is known}
Erd\H{o}s and Gy\'arf\'as~\cite{ErGy99} proved
Conjecture~\ref{conj:main} for $r=3$ (via Brooks' theorem and $R(3,3)=6$) and for
$r=4$ (via $R(3,4)=9$, the uniqueness of its extremal graph, and an intricate
edge count), and observed that it is \emph{false} for $r=2$: the pentagon
colouring of $K_5$, colour $0$ on a $5$-cycle and colour $1$ on the complementary
$5$-cycle, is balanced, so $\Ncap(2)=5=2^2+1$. They also observed a lower bound
that makes the ``$+1$'' essential. From an affine plane of order $r$ one builds a
balanced $r$-colouring of $K_{r^2}$: identify the vertices with the $r^2$ points,
give each of the $r-1$ colours $1,\dots,r-1$ one parallel class of lines, and
give colour $0$ the union of the remaining two parallel classes. Each class is a
disjoint union of cliques (or, for colour $0$, a rook's graph), so has
independence number exactly $r$; hence every $(r+1)$-set sees all $r$ colours.
Affine planes of order $r$ exist for every prime power $r$, so $\Ncap(r)\ge r^2$
for infinitely many $r$. Consequently
\[
  \text{the analogue of Conjecture~\ref{conj:main} with } r^2+1
  \text{ replaced by } r^2 \text{ is false for infinitely many } r,
\]
and any proof of the conjecture must distinguish $r^2$ from $r^2+1$.

For $r=5$ the affine construction gives a balanced $5$-colouring of $K_{25}$, so
$\Ncap(5)\ge 25$, one vertex short of a counterexample. This case is genuinely
harder than $r=3,4$: the matching upper-bound construction of~\cite{ErGy99}
requires a projective plane of order $r+1=6$, which does not exist
(Bruck--Ryser, equivalently the non-existence of two orthogonal Latin squares of
order $6$), so even the finiteness statements available for $r=3,4$ are not
available here. Erd\H{o}s and Gy\'arf\'as remarked that the ``flexibility'' of
the affine construction ``might suggest that the conjecture is not true.''
Problem~$617$ is listed as \textsc{falsifiable} (open, but disprovable by a finite
counterexample) at \cite{Bloom617}, and $r=5$ is its first open case.

\subsection{Results}
Our main result settles the case $r=5$ of Conjecture~\ref{conj:main} on the
impossibility side.

\begin{maintheorem}\label{thm:main}
There is no balanced $5$-colouring of $K_{26}$. Equivalently, $\Ncap(5)=25$: the
affine-plane colourings of $K_{25}$ are optimal, and every $5$-colouring of
$K_{26}$ has $6$ vertices whose induced $K_6$ misses at least one colour.
\end{maintheorem}

Together with~\cite{ErGy99} this establishes Conjecture~\ref{conj:main} for
$r\in\{3,4,5\}$, and confirms $\Ncap(r)=r^2$ for these three values.

Our second result concerns the extremal near-miss. The affine colourings are the
only balanced colourings of $K_{r^2}$ known in the literature, and they come with
a choice --- which two parallel classes to merge --- that a counterexample would
have to exploit. We show they never help.

\begin{maintheorem}\label{thm:affine}
Let $q\ge 3$ be a prime power and let $\mathcal C$ be the balanced $q$-colouring
of $K_{q^2}$ obtained from $\mathrm{AG}(2,q)$ by colouring each edge by the
direction of the line through its endpoints and merging two of the $q+1$
directions into one colour. Then $\mathcal C$ has no balanced one-vertex
extension to $K_{q^2+1}$.
\end{maintheorem}

Theorem~\ref{thm:affine} holds for all prime powers simultaneously, by a short
finite-geometry argument (Section~\ref{sec:affine}); for $r=5$ it says that a
counterexample to Conjecture~\ref{conj:main}, if one existed, could not arise by
adding a vertex to an affine colouring of $K_{25}$. It is strictly weaker than the
$r=5$ case of the conjecture --- it does not rule out non-affine balanced
colourings of $K_{25}$ --- but it is clean, uniform in $q$, and independent of the
machinery behind Theorem~\ref{thm:main}.

\subsection{Why \texorpdfstring{$r=5$}{r=5} resists the earlier method}\label{ssec:why}
Erd\H{o}s and Gy\'arf\'as settled $r=3,4$ by attacking a \emph{minority} colour
class of $K_{r^2+1}$ (one with the fewest edges) directly and forcing an
independent $(r+1)$-set inside it by hand: for $r=3$ from Brooks' theorem and
$R(3,3)=6$, and for $r=4$ from $R(3,4)=9$, the uniqueness of its extremal graph,
and an intricate edge count. Neither ingredient scales cleanly to $r=5$. That
analysis is organised around Ramsey numbers with unique extremal graphs, and at
$r=5$ the relevant inputs multiply --- $R(3,5)$, $R(4,4)$, $R(4,5)$, and more ---
with no single extremal structure to anchor them; and, as noted above, the
matching upper-bound construction needs a projective plane of order $r+1=6$, which
does not exist, so even the finiteness statement available for $r=3,4$ is out of
reach here.

What makes the case go through is a change of object. The reduction of
Section~\ref{sec:reduction} deletes a vertex and folds the surrounding colouring
into two self-contained statements about a \emph{single} graph on $25$ vertices
(Lemmas~\ref{lem:MH} and~\ref{lem:MM}); thereafter one classical extremal input ---
Brouwer's non-partite Tur\'an bound with the Kang--Pikhurko equality case ---
drives a uniform counting recursion in place of the earlier case-by-case analysis,
and only four facts about graphs on at most $12$ vertices are left to a SAT solver
(Section~\ref{ssec:method} sets out the layers). The quantitative thresholds in the
two lemmas were not guessed but read off a computation: they sit exactly at an
empirical phase transition (Appendix~\ref{app:phase}). The ``hitter'' configurations
that Lemma~\ref{lem:MH} forbids exist for every $n\le 24$ and disappear at precisely
$n=25$ --- the conjecture's ``$+1$'' costing one vertex of slack --- and the same
``$+1$'' resurfaces in the proof as the arithmetic $21=4\cdot 5+1$, forcing one
oversized part where $20=4\cdot 5$ splits evenly (Section~\ref{ssec:MH-n24}). This
is where the argument uses the ``$+1$'' essentially, as any proof must: with the
colouring forgotten, both lemmas become false (Appendix~\ref{app:phase}).

\subsection{Method and provenance}\label{ssec:method}
The proof of Theorem~\ref{thm:main} has three layers, described in outline here
and in detail below.

\emph{The reduction} (Section~\ref{sec:reduction}). Suppose $K_{26}$ had a
balanced $5$-colouring. Delete a vertex $x$; the induced colouring of $K_{25}$ is
balanced, and partitioning the $25$ remaining vertices by the colour of their edge
to $x$ produces five sets $T_0,\dots,T_4$. An elementary argument
(Lemma~\ref{lem:chain}) shows that two structural statements about the colour
classes on $25$ vertices --- Lemmas~\ref{lem:MH} (``\textsf{[MH$''$]}'') and
\ref{lem:MM} (``\textsf{[MM]}'') below --- force a contradiction. This step is short and
entirely elementary.

\emph{The two lemmas} (Sections~\ref{sec:MH} and~\ref{sec:MM}). Each is a
statement about graphs on $\le 25$ vertices in which every $6$-set spans at most
$11$ edges (we call this \capeleven), with prescribed independence numbers. Their
proofs share a double-counting recursion (Section~\ref{sec:toolkit}) built from
\capeleven{} and one classical theorem: Brouwer's bound on the number of edges in a
$K_{r+1}$-free graph that is not $r$-partite, together with the Kang--Pikhurko
classification of the extremal graphs~\cite{Brouwer81,KP05}. The recursion bottoms
out in four finite facts about small graphs, which we verify by SAT.

\emph{The certificates and formalisation} (Sections~\ref{sec:certificates}
and~\ref{sec:lean}). The four finite facts are certified by DRAT/LRAT proofs,
independently checked. The whole argument --- reduction, both lemmas, the
recursion, the SAT facts (through Lean's verified certificate checker), Brouwer's
edge bound, and the Kang--Pikhurko equality classification at $(5,21)$ --- has been
formalised in Lean~4. The headline theorem \texttt{erdos\_617\_r5\_unconditional} is
\texttt{sorry}-free and carries no mathematical hypothesis; it is unconditional in the
exact sense that it depends only on the three standard Lean axioms and the
\texttt{native\_decide} reflection axioms (four for the SAT certificates, ten for the
finite Kang--Pikhurko construction facts). It is thus a machine-checked proof modulo
that disclosed trust base, not an appeal to any unformalised input; a conditional form
taking the classification as an explicit hypothesis is retained for readers who prefer
to audit it against the literature (Section~\ref{ssec:lean-brouwer}).

\medskip
This document is a \emph{verification-ready report}, not an authored announcement,
and its whole purpose is to be checked by others. The mathematics was produced and
internally reviewed by AI systems, which we disclose in full in
Section~\ref{sec:provenance} in keeping with the norms of \cite{Bloom617}. In brief:
the proofs of the two lemmas (Sections~\ref{sec:MH}--\ref{sec:MM}) were drafted by a
large language model (OpenAI \texttt{gpt-5.6-sol}) from self-contained briefs; the
reduction, the verification and review infrastructure, the SAT certification, and the
Lean formalisation were produced by Claude (Anthropic) sessions; and each of the
three links was subjected to an independent adversarial review by a fresh model
session, whose findings we incorporate (both lemmas required a repair, and those
repairs are what we present). This internal chain --- author sessions, adversarial
reviewers, exact recomputation, and a Lean formalisation --- is a strong but
\emph{not yet independently human-refereed} certificate. We make no claim that it
substitutes for expert scrutiny; the document is offered precisely so that such
scrutiny is easy to perform. Section~\ref{sec:howto} is a verification guide for the
reader who wants to check the claim, at several depths.

\subsection{Organisation}
Section~\ref{sec:howto} is a verification guide: how to check the claim in $30$
minutes (machine), in $120$ minutes (by reading), or exhaustively.
Section~\ref{sec:reduction} fixes notation, records the elementary properties of
balanced colourings, states the two lemmas, and proves Theorem~\ref{thm:main} from
them. Section~\ref{sec:toolkit} develops the shared \capeleven{} toolkit
(Brouwer's theorem, the counting identity, and the $M$- and $L$-tables).
Sections~\ref{sec:MH} and~\ref{sec:MM} prove Lemmas~\ref{lem:MH} and~\ref{lem:MM}.
Section~\ref{sec:affine} proves Theorem~\ref{thm:affine}.
Section~\ref{sec:certificates} documents the four SAT certificates and the
recipe to regenerate and re-check them. Section~\ref{sec:lean} describes the Lean
formalisation and its axiom profile. Section~\ref{sec:provenance} gives the full
provenance and an honest account of the argument's verification status.
Appendix~\ref{app:phase} records the empirical phase-transition data that
motivated the shape of the argument.

%% ===========================================================================
\section{How to verify this report}\label{sec:howto}
%% ===========================================================================

Because this is a report offered for checking rather than a claim asking to be
trusted, we describe up front how a reader can verify it, at three depths. The
accompanying repository contains the Lean project, the SAT certificates, and the
arithmetic scripts referenced below.

\subsection*{The $30$-minute path (statement audit plus machine check)}
This path certifies the theorem \emph{modulo} the two clearly delimited
non-elementary inputs, without reading the combinatorial proofs.
\begin{enumerate}[label=\textup{(\arabic*)},leftmargin=2.2em]
\item \emph{Audit the statement.} Read Theorem~\ref{thm:main} and the Lean statements
\texttt{Main} and \texttt{erdos\_617\_r5\_\allowbreak upstream} (Section~\ref{sec:lean}), and
confirm they say ``every $5$-colouring of $K_{26}$ has a $6$-set missing a colour'' ---
matching the \texttt{erdos\_617} statement of the \texttt{formal-conjectures}
repository at $r=5$~\cite{FC}. A wrong statement would invalidate everything, so this
step comes first.
\item \emph{Build and check axioms.} In \texttt{lean617/}, run \texttt{lake exe cache
get \&\& lake build} (a clean build), then \texttt{\#print axioms
Erdos617.erdos\_617\_r5\_\allowbreak unconditional}. The result should be
\texttt{propext}, \texttt{Classical.choice}, \texttt{Quot.sound}, plus fourteen named
\texttt{native\_decide} reflection axioms (four for the SAT certificates and ten for
the finite Kang--Pikhurko construction facts) --- \emph{seventeen} in all --- with
\emph{no} mathematical hypothesis in the theorem's type and, crucially, \emph{no}
\texttt{sorryAx}. (The still-exported conditional \texttt{erdos\_617\_r5} instead
carries \texttt{KPEqualityClassification} as an explicit hypothesis and only the three
standard plus four SAT axioms, for a reader who prefers to supply the classification
themselves.) The absence of \texttt{sorryAx} is the definitive ``no admitted proof''
check; a textual grep for \texttt{sorry} finds the token only inside
\texttt{sorry-free} annotations in the sources.
\end{enumerate}
What this establishes, stated exactly: \emph{the headline theorem
\texttt{erdos\_617\_r5\_unconditional} is \texttt{sorry}-free and carries no
mathematical hypothesis, resting only on the three standard Lean axioms and fourteen
\texttt{native\_decide} reflection axioms (seventeen in all).} It is
``unconditional'' strictly in that sense --- no unformalised mathematical input ---
and not a claim of a kernel-only formal proof: the four SAT certificates and the ten
finite Kang--Pikhurko construction facts are trusted through \texttt{native\_decide}'s
compiler-level reflection (Section~\ref{sec:lean}). Both Brouwer's edge bound and the
equality classification, formerly assumed, are now proved in Lean. What remains for a
reader after this path is exactly to satisfy themselves of those
\texttt{native\_decide} inputs --- the next two paths do that.

\subsection*{The $120$-minute path (read the mathematics)}
The proof is elementary and self-contained apart from one classical theorem. Read
Section~\ref{sec:reduction} (the reduction; short), then the two lemma proofs of
Sections~\ref{sec:MH} and~\ref{sec:MM}, using Section~\ref{sec:toolkit} as needed.
The only external input is Brouwer's theorem (Theorem~\ref{thm:brouwer}), a published
result whose five uses are itemised in Section~\ref{sec:MH}; the only computational
inputs are the four finite facts of Fact~\ref{fact:sat}, about graphs on at most $12$
vertices. Everything else is by hand and is confirmed by exact recomputation
(\texttt{tools/verify\_gpt\_arith.py}). A reader who follows this path and the
statement audit above has checked the mathematics end to end, modulo trusting (a) the
cited theorem and (b) the four SAT certificates --- both discharged next.

\subsection*{The completist path (audit the machine inputs)}
Two things are trusted by the shorter paths and can be checked directly.
\emph{The SAT facts:} regenerate the four DRAT certificates and check them with
\texttt{drat-trim} (Section~\ref{sec:certificates}, path~(i)); this is independent of
Lean and of \texttt{native\_decide}. \emph{The classical input:} the informal proof
uses Brouwer's bound and the Kang--Pikhurko equality classification, both from
\cite{Brouwer81,KP05} (the repository records this check use-by-use); in the Lean
development both are now \emph{proved} rather than assumed
(Section~\ref{ssec:lean-brouwer}), so a reader who trusts that formalisation need audit
neither, while a reader who would rather lean on the literature can check the published
classification and consume it through the conditional theorem instead. A reader who
completes all three paths has reduced the entire claim to the correctness of
\texttt{drat-trim}, the Lean kernel, and \texttt{native\_decide}'s compiler-level
reflection on the fourteen fixed certificates and construction facts.

\subsection*{What is deliberately \emph{not} yet claimed}
This report does \emph{not} claim the result is community-accepted or independently
refereed by a human expert, nor that the Lean proof is free of a trust base: it is
unconditional only in the sense that no \emph{mathematical} hypothesis remains, while
the four SAT certificates and the ten finite Kang--Pikhurko construction facts are
still trusted through \texttt{native\_decide}'s compiler-level reflection
(Section~\ref{sec:lean}). The Kang--Pikhurko equality classification at $(5,21)$,
formerly carried as the hypothesis \texttt{KPEqualityClassification}, is now discharged
in Lean (Section~\ref{ssec:lean-brouwer}) --- but that discharge is a fresh,
load-bearing formalisation authored by the same AI pipeline as the rest and reviewed
only internally (Section~\ref{sec:provenance}). Those two --- the internal-only review
and the \texttt{native\_decide} trust base --- are the gaps a third-party verifier is
invited to close or to press on.

%% ===========================================================================
\section{Definitions and the reduction}\label{sec:reduction}
%% ===========================================================================

\subsection{Notation and conventions}
All graphs are finite and simple. For a graph $G$ we write $V(G)$, $e(G)=\abs{E(G)}$,
$d_G(v)$ (or $d_v$) for the degree of $v$, $N_G(v)$ for its neighbourhood,
$N_G[v]=N_G(v)\cup\{v\}$, $\alpha(G)$ for the independence number, and $\omega(G)$
for the clique number. For $S\subseteq V(G)$, $G[S]$ is the induced subgraph and
$e_G(S)=e(G[S])$; for disjoint $A,B$, $e_G(A,B)$ is the number of edges with one
endpoint in each. We write $t_r(n)=\ex(n,K_{r+1})$ for the Tur\'an number (the
maximum number of edges in a $K_{r+1}$-free graph on $n$ vertices) and $T_r(n)$ for
the Tur\'an graph. All colourings are edge-colourings of a complete graph, and
colours are labelled $0,\dots,r-1$.

For a $5$-colouring $\chi$ we write $G_0,\dots,G_4$ for the colour classes. Two
properties recur so often that we name them.

\begin{definition}[\capeleven]
A graph $G$ (typically a colour class of a colouring of $K_n$) is \emph{\capeleven}
if every $6$-vertex subset spans at most $11$ edges of $G$.
\end{definition}

\begin{definition}[hitter]
Let $G$ be a graph and $k\ge 0$ an integer. A vertex set $T$ is a \emph{$k$-hitter}
for $G$ if $\alpha(G-T)\le 4$ and $\abs{T}=k$; that is, deleting the $k$ vertices
of $T$ leaves a graph with no independent $5$-set. (The value $4$ is fixed
throughout, tuned to $r=5$.)
\end{definition}

\subsection{Elementary properties of balanced colourings}
The following facts are immediate from the definition of balancedness and are used
constantly.

\begin{lemma}\label{lem:elem}
Let $\chi$ be a balanced $5$-colouring of $K_n$ with colour classes
$G_0,\dots,G_4$. Then:
\begin{enumerate}[label=\textup{(\alph*)}]
\item\label{it:restrict} the restriction of $\chi$ to any subset of vertices is a
balanced $5$-colouring of the corresponding complete graph;
\item\label{it:alpha5} $\alpha(G_c)\le 5$ for every colour $c$;
\item\label{it:cap} every $G_c$ is \capeleven{} and $K_6$-free;
\item\label{it:minority} if $n=25$ then $\sum_{c}e(G_c)=\binom{25}{2}=300$, so
some colour $m$ (a \emph{minority} colour) has $e(G_m)\le 60$.
\end{enumerate}
\end{lemma}

\begin{proof}
\ref{it:restrict} A subset of an $(r+1)$-set of the smaller graph is an
$(r+1)$-set of $K_n$; its edge colours are unchanged, so it still sees all
$5$ colours.

\ref{it:alpha5} An independent $6$-set of $G_c$ is a $6$-set spanning no colour-$c$
edge, hence missing colour $c$ --- impossible.

\ref{it:cap} A $6$-set has $\binom 62=15$ edges. If $12$ of them had colour $c$,
the remaining $3$ could realise at most $3$ further colours, so the $6$-set would
see at most $4$ colours; hence at most $11$ edges of any one colour, which is
\capeleven. A monochromatic $K_6$ sees a single colour, so no class contains one.

\ref{it:minority} The five classes partition the $\binom{25}{2}=300$ edges, so the
minimum class size is at most the average $300/5=60$. (Ties are broken
arbitrarily.)
\end{proof}

Note that \ref{it:cap} makes ``\capeleven'' and ``$K_6$-free'' standing hypotheses
on \emph{every} class of a balanced $5$-colouring, and they are inherited by every
induced subgraph.

\subsection{The two lemmas}
The reduction turns Theorem~\ref{thm:main} into two statements that no longer
mention the deleted vertex. The first forbids a small hitter in a single class of a
balanced colouring of $K_{25}$.

\begin{lemma}[\textsf{[MH$''$]}]\label{lem:MH}
No balanced $5$-colouring of $K_{25}$ has a colour $c$ and a $4$-hitter $T$ for
$G_c$. Equivalently: there is no $5$-colouring of $K_{21}$ in which every class is
\capeleven{} and $K_6$-free and the colour classes have independence numbers
$(4,5,5,5,5)$.
\end{lemma}

The second is a purely graph-theoretic statement about the minority class, phrased
so as to forget the surrounding colouring; the reduction supplies each of its
hypotheses.

\begin{lemma}[\textsf{[MM]}]\label{lem:MM}
There is no graph $G$ on $25$ vertices with $\alpha(G)\le 5$, \capeleven, and
$e(G)\le 60$, together with a $5$-vertex set $T$ such that $\alpha(G-T)\le 4$ and
$e(G[T])\le 6$.
\end{lemma}

Lemma~\ref{lem:MH} is proved in Section~\ref{sec:MH} and Lemma~\ref{lem:MM} in
Section~\ref{sec:MM}. We first show they suffice.

\subsection{Proof of Theorem~\ref{thm:main}}
\begin{lemma}[Reduction]\label{lem:chain}
If Lemmas~\ref{lem:MH} and~\ref{lem:MM} hold, then $K_{26}$ has no balanced
$5$-colouring.
\end{lemma}

\begin{proof}
Suppose $\chi$ is a balanced $5$-colouring of $K_{26}$. Fix a vertex $x$, let
$V'=V(K_{26})\setminus\{x\}$ (so $\abs{V'}=25$), and let $\chi'=\chi|_{V'}$, a
balanced $5$-colouring of $K_{25}$ by Lemma~\ref{lem:elem}\ref{it:restrict}. Let
$G_0,\dots,G_4$ be the colour classes of $\chi'$ and, for each colour $c$, let
\[
  T_c \;=\; \{\, v\in V' : \chi(xv)=c \,\}.
\]
The sets $T_0,\dots,T_4$ partition $V'$, so $\sum_c\abs{T_c}=25$.

\smallskip
\noindent\emph{Step 1: $\alpha(G_c-T_c)\le 4$ for every $c$.}
If $F$ were an independent $5$-set of $G_c$ with $F\cap T_c=\varnothing$, then the
$6$-set $\{x\}\cup F$ would miss colour $c$ under $\chi$: no edge inside $F$ has
colour $c$ (as $F$ is independent in $G_c$), and no edge $xv$ with $v\in F$ has
colour $c$ (as $v\notin T_c$). This contradicts balancedness of $\chi$. So $T_c$ is
a hitter for $G_c$ in the sense that $\alpha(G_c-T_c)\le 4$. This is the one place
the argument uses the full $K_{26}$.

\smallskip
\noindent\emph{Step 2: $\abs{T_c}=5$ for every $c$.}
Suppose some $\abs{T_c}\le 4$. Choose a $4$-set $T^+\supseteq T_c$ inside $V'$
(possible, as $\abs{V'}=25$). Deleting extra vertices cannot raise an independence
number, so $\alpha(G_c-T^+)\le\alpha(G_c-T_c)\le 4$ by Step~1; thus $T^+$ is a
$4$-hitter for $G_c$ in the balanced colouring $\chi'$, contradicting
Lemma~\ref{lem:MH}. Hence $\abs{T_c}\ge 5$ for all $c$, and since the five sizes sum
to $25$, all equal $5$.

\smallskip
\noindent\emph{Step 3: each $T_{c'}$ spans at most $6$ edges of colour $c'$.}
Fix distinct colours $c\ne c'$. The set $T_{c'}$ has size $5$ and is disjoint from
$T_c$. If $T_{c'}$ contained no colour-$c$ edge, it would be an independent $5$-set
of $G_c$ disjoint from $T_c$, contradicting $\alpha(G_c-T_c)\le 4$ (Step~1). Hence
$T_{c'}$ contains at least one edge of colour $c$, for each of the four colours
$c\ne c'$. Edges of distinct colours are distinct, so at least $4$ of the
$\binom 52=10$ edges inside $T_{c'}$ have a colour other than $c'$, leaving at most
$6$ edges of colour $c'$.

\smallskip
\noindent\emph{Step 4: contradiction.}
Let $m$ be a minority colour, so $e(G_m)\le 60$
(Lemma~\ref{lem:elem}\ref{it:minority}). By
Lemma~\ref{lem:elem}\ref{it:alpha5}--\ref{it:cap}, $G_m$ has $\alpha(G_m)\le 5$ and
is \capeleven. By Step~2, $\abs{T_m}=5$; by Step~1, $\alpha(G_m-T_m)\le 4$; and by
Step~3 (with $c'=m$), $e(G_m[T_m])\le 6$. Thus $(G_m,T_m)$ satisfies every
hypothesis of Lemma~\ref{lem:MM}, which forbids exactly this configuration.
Contradiction.
\end{proof}

Theorem~\ref{thm:main} now follows: Lemmas~\ref{lem:MH} and~\ref{lem:MM} are proved
below, so Lemma~\ref{lem:chain} applies. The equality $\Ncap(5)=25$ combines this
with the affine construction ($\Ncap(5)\ge 25$) and monotonicity.

\begin{remark}
Every hypothesis of Lemma~\ref{lem:MM} is used: $e(G_m)\le 60$ is the minority
bound, $\abs{T_m}=5$ comes from Lemma~\ref{lem:MH} via Step~2 (a $4$-point $T_m$
would contain no $5$-set and Step~3 would be vacuous), and $e(G_m[T_m])\le 6$ is
the ``usability'' bound of Step~3. The pure-graph relaxations of both lemmas ---
dropping the surrounding colouring --- are \emph{false}: there is an $80$-edge graph
on $21$ vertices with $\alpha\le 4$ and \capeleven{} (so Lemma~\ref{lem:MH} genuinely
needs balancedness), and $(\alpha\le 5,\ \capeleven,\ \le 60\text{ edges})$ alone is
satisfiable on $25$ vertices (so Lemma~\ref{lem:MM} genuinely needs the hitter and
usability hypotheses). See Appendix~\ref{app:phase}.
\end{remark}

%% ===========================================================================
\section{The \texorpdfstring{\capeleven}{cap-11} toolkit}\label{sec:toolkit}
%% ===========================================================================

Both lemmas are proved by bounding the number of edges of a \capeleven{} graph
below in terms of its independence number and clique number. This section
assembles the shared machinery: one classical extremal theorem, a double-counting
identity, a neighbourhood bound coming from \capeleven, four finite facts about
very small graphs, and two tables of edge lower bounds ($M$ for independence number
$2$ and $L$ for independence number $3$) that the two lemmas consume. Throughout,
``\capeleven'' means every $6$-set spans at most $11$ edges, a property inherited
by all induced subgraphs.

\subsection{Brouwer's non-partite Tur\'an theorem}
A $K_{r+1}$-free graph on $n$ vertices has at most $t_r(n)$ edges, with equality
only for the Tur\'an graph $T_r(n)$, which is $r$-partite. If one additionally
forbids $r$-partiteness, the bound drops.

\begin{theorem}[Brouwer~\cite{Brouwer81}; equality by Kang--Pikhurko~\cite{KP05}]
\label{thm:brouwer}
Let $r\ge 2$ and $n\ge 2r+1$. If $G$ is a $K_{r+1}$-free graph on $n$ vertices
that is not $r$-partite (equivalently $\chi(G)>r$), then
\[
  e(G)\;\le\; t_r(n)-\Big\lfloor\tfrac nr\Big\rfloor+1 .
\]
Moreover, Kang and Pikhurko classify the extremal graphs: every graph attaining
equality arises from an explicit construction $G(\mathbf n)$ obtained from a
complete $r$-partite graph on parts $N_1,\dots,N_r$ (with $\sum\abs{N_i}=n-1$) by
adding one vertex and deleting a matching-like set of edges, with the part sizes
$\mathbf n$ ranging over an explicit short list.
\end{theorem}

We use Theorem~\ref{thm:brouwer} only for $r=5$ and $n\in\{15,16,21\}$, all of which
satisfy $n\ge 2r+1=11$.\footnote{Brouwer's $1981$ report~\cite{Brouwer81} is a
Mathematisch Centrum technical note and is not readily available online; we quote
the bound in the form in which it is stated and attributed to Brouwer by Ren, Wang,
Wang, and Yang~\cite{RWWY}, and take the equality classification from the primary
source~\cite{KP05}. All Tur\'an arithmetic is re-derived in the repository.} The
relevant Tur\'an numbers are
\[
  t_5(15)=90,\qquad t_5(16)=102,\qquad t_5(21)=176,
\]
computed from the balanced part sizes $(3,3,3,3,3)$, $(4,3,3,3,3)$ and
$(5,4,4,4,4)$ respectively, giving Brouwer bounds
\begin{equation}\label{eq:brouwer-values}
  90-3+1=88,\qquad 102-3+1=100,\qquad 176-4+1=173.
\end{equation}
We use the equality classification only once, at $(r,n)=(5,21)$, in the following
form; it is a direct reading of the Kang--Pikhurko construction and is verified
against their Lemma~5 in Section~\ref{sec:MH} where it is applied.

\begin{corollary}[$n=21$ equality shape]\label{cor:equality21}
Let $F$ be a $K_5$-free graph on $21$ vertices with $\alpha(F)\le 5$ and exactly
$e(F)=37$ edges. Then there are disjoint sets $A,B$ with $\abs A=5$, $\abs B=4$
such that $F[A]=K_5-xy$ for some $x,y\in A$, $F[B]=K_4$, and exactly $4$ edges of
$F$ join $A$ to $B$.
\end{corollary}

\subsection{A counting identity}
For a graph $G$ on $q$ vertices and a vertex $v$, write
\[
  W_v \;=\; V(G)\setminus N_G[v]
\]
for the set of non-neighbours of $v$ other than $v$ itself, so $\abs{W_v}=q-1-d_v$.
The following two identities are standard double counts; $\tau(G)$ denotes the
number of triangles.

\begin{lemma}\label{lem:identity}
For every graph $G$ on $q$ vertices,
\begin{align}
  \sum_{v} e(G[W_v]) &= q\,e(G)-\sum_v d_v^2+3\tau(G),\label{eq:41}\\
  3\tau(G) &= \sum_v e(G[N(v)]).\label{eq:42}
\end{align}
\end{lemma}

\begin{proof}
For~\eqref{eq:41}: a fixed edge $xy$ lies in $G[W_v]$ exactly when $v$ is
non-adjacent to both $x$ and $y$ and $v\notin\{x,y\}$, i.e.\ for
$q-\abs{N(x)\cup N(y)}=q-d_x-d_y+\abs{N(x)\cap N(y)}$ vertices $v$ (using
$\{x,y\}\subseteq N(x)\cup N(y)$). Summing over edges $xy$ gives
$\sum_v e(G[W_v]) = q\,e(G)-\sum_{xy\in E}(d_x+d_y)+\sum_{xy\in E}\abs{N(x)\cap N(y)}$.
The middle sum is $\sum_v d_v^2$, and $\abs{N(x)\cap N(y)}$ counts triangles on the
edge $xy$, so the last sum is $3\tau(G)$. For~\eqref{eq:42}: each triangle is
counted once for each of its three vertices $v$, as an edge of $G[N(v)]$.
\end{proof}

\subsection{The neighbourhood bound from \texorpdfstring{\capeleven}{cap-11}}
\capeleven{} limits how many edges a neighbourhood can carry.

\begin{lemma}\label{lem:nbhd}
Let $G$ be \capeleven{} and $v$ a vertex of degree $d$. Then
\[
  e(G[N(v)])\;\le\; b(d)\;:=\;
  \begin{cases}
    \binom d2, & d\le 4,\\[1mm]
    \big\lfloor \tfrac{3d(d-1)}{10}\big\rfloor, & d\ge 5.
  \end{cases}
\]
If in addition $\omega(G)\le 4$, then $G[N(v)]$ is $K_4$-free, so we may replace
$b(d)$ by the smaller bound $u(d):=\min\{b(d),\,\ex(d,K_4)\}$, where $\ex(d,K_4)$ is
the Tur\'an number.
\end{lemma}

\begin{proof}
For $d\le 4$ the bound $\binom d2$ is trivial. For $d\ge 5$, let $A\subseteq N(v)$
be any $5$-set. The $6$-set $\{v\}\cup A$ carries $5$ edges from $v$ to $A$ and
$e(G[A])$ edges inside $A$, so \capeleven{} gives $5+e(G[A])\le 11$, i.e.\
$e(G[A])\le 6$. Summing over all $\binom d5$ five-subsets of $N(v)$, each edge of
$G[N(v)]$ lying in $\binom{d-2}{3}$ of them,
\[
  e(G[N(v)])\binom{d-2}{3}\;\le\; 6\binom d5,
\]
and $6\binom d5/\binom{d-2}3=\tfrac{3d(d-1)}{10}$ gives the bound. If $\omega(G)\le
4$ then $G[N(v)]$ has no $K_4$ (a $K_4$ in $N(v)$ plus $v$ is a $K_5$), so
$e(G[N(v)])\le\ex(d,K_4)$ as well.
\end{proof}

\subsection{Four certified small-graph facts}\label{ssec:small}
The recursion below rests on four facts about small \capeleven{} graphs with
independence number $2$. Each is a finite statement; each is proved by SAT and
independently DRAT-checked (Section~\ref{sec:certificates}).

\begin{fact}\label{fact:sat}
\leavevmode
\begin{enumerate}[label=\textup{(M\arabic*)}]
\item\label{sat:M9} Every \capeleven{} graph on $9$ vertices with $\alpha\le 2$,
$\omega\le 4$ has at least $19$ edges.
\item\label{sat:M10} Every \capeleven{} graph on $10$ vertices with $\alpha\le 2$,
$\omega\le 4$ has at least $25$ edges.
\item\label{sat:nonex11} There is no \capeleven{} graph on $11$ vertices with
$\alpha\le 2$ \textup{(no clique bound needed)}.
\item\label{sat:nonex12} There is no \capeleven{} graph on $12$ vertices with
$\alpha\le 2$ \textup{(no clique bound needed)}.
\end{enumerate}
\end{fact}

\begin{remark}\label{rem:nonex-mono}
Facts~\ref{sat:nonex11} and~\ref{sat:nonex12} are $\omega$-free: they hold for any
clique number. Since $\alpha\le 2$ and \capeleven{} are both inherited by induced
subgraphs, \ref{sat:nonex11} alone implies that \emph{no} \capeleven{} graph on
$n\ge 11$ vertices has $\alpha\le 2$; \ref{sat:nonex12} is a redundant independent
check. This $\omega$-freeness is what the repaired proofs of both lemmas use
(Sections~\ref{ssec:MH-K5free} and~\ref{ssec:MM-two}); the clique bound $\omega\le
4$ is genuinely required in \ref{sat:M9} and~\ref{sat:M10}, where it is available
in context, since $K_5\sqcup K_4$ and $K_5\sqcup K_5$ are \capeleven{} graphs with
$\alpha=2$ on $9$ and $10$ vertices.
\end{remark}

\subsection{The $M$-table and the $L$-table}\label{ssec:tables}
We package the edge lower bounds into two functions.

\begin{definition}
For $s\ge 0$ let
\[
  M(s)\;=\;\min\{\, e(Z) : \abs{V(Z)}=s,\ \alpha(Z)\le 2,\ \omega(Z)\le 4,\
  Z\text{ \capeleven}\,\},
\]
with $M(s)=+\infty$ if no such graph exists, and let $L(s)$ be defined the same
way with $\alpha\le 3$ in place of $\alpha\le 2$.
\end{definition}

\begin{lemma}\label{lem:Mtable}
$M(s)=\binom s2-\lfloor s^2/4\rfloor$ for $0\le s\le 8$; $M(9)=19$; $M(10)=25$;
and $M(s)=+\infty$ for $s\ge 11$.
\end{lemma}

\begin{proof}
For $s\le 8$, $\alpha(Z)\le 2$ means the complement $\overline Z$ is triangle-free,
so by Mantel's theorem $e(\overline Z)\le\lfloor s^2/4\rfloor$ and
$e(Z)\ge\binom s2-\lfloor s^2/4\rfloor$; this bound is attained (e.g.\ by the
complement of a balanced complete bipartite graph, which is \capeleven{} and
$K_4$-free for $s\le 8$), and \capeleven{} is vacuous in this range. The values
$M(9)=19$ and $M(10)=25$ are Facts~\ref{sat:M9} and~\ref{sat:M10} (attainment is by
explicit graphs). Finally $M(s)=+\infty$ for $s\ge 11$ is
Fact~\ref{sat:nonex11} and Remark~\ref{rem:nonex-mono}.

For completeness we record the $s\le 10$ values, which are all we use:
\[
  \begin{array}{c|cccccccccc}
    s & 1 & 2 & 3 & 4 & 5 & 6 & 7 & 8 & 9 & 10\\\hline
    M(s) & 0 & 0 & 1 & 2 & 4 & 6 & 9 & 12 & 19 & 25.
  \end{array}\qedhere
\]
\end{proof}

\begin{lemma}\label{lem:Ltable}
For $s\le 12$, $L(s)=\binom s2-t_3(s)$ \textup{(the complement-Tur\'an value for
$\alpha\le 3$)}. For $13\le s\le 20$,
\[
  \begin{array}{c|cccccccc}
    s & 13 & 14 & 15 & 16 & 17 & 18 & 19 & 20\\\hline
    L(s) & 24 & 31 & 38 & 46 & 53 & 62 & 73 & 84 .
  \end{array}
\]
\end{lemma}

\begin{proof}
For $s\le 12$, $\alpha\le 3$ means $\overline Z$ is $K_4$-free, so
$e(\overline Z)\le t_3(s)=\ex(s,K_4)$ and $e(Z)\ge\binom s2-t_3(s)$, attained by the
complement of the Tur\'an graph $T_3(s)$ (which is a disjoint union of $\le 3$
cliques, so \capeleven{} and $K_4$-free once $s\le 12$). This gives the small values
\[
  \begin{array}{c|cccccccccccc}
    s & 3 & 4 & 5 & 6 & 7 & 8 & 9 & 10 & 11 & 12\\\hline
    L(s) & 0 & 1 & 2 & 3 & 5 & 7 & 9 & 12 & 15 & 18 .
  \end{array}
\]

For $13\le s\le 20$ we run the recursion of Lemma~\ref{lem:identity}. Let $X$ be a
\capeleven{} graph on $s$ vertices with $\alpha(X)\le 3$ and $\omega(X)\le 4$. For
each $v$, $W_v$ has $\alpha(W_v)\le 2$ (an independent $3$-set in $W_v$ plus $v$
would be an independent $4$-set), $\omega(W_v)\le 4$, and is \capeleven; so
$e(X[W_v])\ge M(s-1-d_v)$, where by Lemma~\ref{lem:Mtable} the value is finite only
for $s-1-d_v\le 10$, i.e.\ $d_v\ge s-11$. Combining~\eqref{eq:41},~\eqref{eq:42} and
Lemma~\ref{lem:nbhd} (in its $\omega\le 4$ form),
\[
  \sum_v M(s-1-d_v)\;\le\;\sum_v e(X[W_v])
  = s\,e(X)-\sum_v d_v^2+\sum_v e(X[N(v)])
  \;\le\; s\,e(X)-\sum_v d_v^2+\sum_v u(d_v),
\]
which rearranges to $\sum_v\Phi_s(d_v)\le 0$, where
\[
  \Phi_s(d)\;=\;M(s-1-d)+d^2-\tfrac s2\,d-u(d)
\]
(and $\Phi_s(d)=+\infty$ when $d<s-11$, forbidding those degrees). Direct
substitution gives affine lower bounds $\Phi_s(d)\ge A_s-B_s d$ valid at every
feasible integer $d$, with
\[
  \begin{array}{c|cccccccc}
    s & 13 & 14 & 15 & 16 & 17 & 18 & 19 & 20\\\hline
    A_s & 25 & 26 & 30 & 34 & 28 & 34 & 57 & 125/2\\
    B_s & 7 & 6 & 6 & 6 & 9/2 & 5 & 15/2 & 15/2 .
  \end{array}
\]
Then $0\ge\sum_v\Phi_s(d_v)\ge sA_s-B_s\sum_v d_v=sA_s-2B_s\,e(X)$, so
$e(X)\ge\lceil sA_s/(2B_s)\rceil$, which equals the tabulated $L(s)$ in each case.
The affine bounds and the ceilings are verified by exact integer arithmetic; see
Section~\ref{sec:certificates} and the script \texttt{tools/verify\_gpt\_arith.py}.
\end{proof}

\begin{remark}
The refinement $u(d)=\min\{b(d),\ex(d,K_4)\}$ (rather than $b(d)$ alone) is needed
at exactly $d=4$, where $u(4)=\ex(4,K_4)=5<6=\binom 42$; elsewhere $u=b$. Using
$b(4)=6$ would break the affine bound at $s\in\{13,14,15\}$. The nonexistence at
$11$ vertices (Fact~\ref{sat:nonex11}) is used only to declare degrees $d<s-11$
infeasible, tightening the feasible range; it makes the floor larger, never
smaller.
\end{remark}

Finally we record the one \capeleven{} consequence used repeatedly to control how
$K_5$'s attach to the rest of a graph.

\begin{lemma}\label{lem:K5indeg}
If $G$ is \capeleven{} and $Q\subseteq V(G)$ induces a $K_5$, then every vertex
$v\notin Q$ has at most one neighbour in $Q$. Consequently, between two disjoint
$K_5$'s the edges form a matching.
\end{lemma}

\begin{proof}
If $v$ had two neighbours in $Q$, the $6$-set $\{v\}\cup Q$ would carry the
$\binom 52=10$ edges of $Q$ plus at least $2$ more, exceeding $11$.
\end{proof}

%% ===========================================================================
\section{Proof of Lemma \textsf{[MH\texorpdfstring{$''$}{''}]}}\label{sec:MH}
%% ===========================================================================

We prove the equivalent statement in Lemma~\ref{lem:MH}: there is no
$5$-colouring of $K_{21}$ whose colour classes are all \capeleven{} and $K_6$-free
and have independence numbers $(4,5,5,5,5)$. (The equivalence is the reduction to
$21$ vertices in Section~\ref{ssec:MH-reduce}; a $4$-hitter $T$ in a balanced
$K_{25}$ removes $4$ vertices, and the remaining $21$ carry exactly such a
colouring by Lemma~\ref{lem:elem}.)

Suppose such a colouring of $K_{21}$ existed. Write $H$ for the class of
independence number $4$ and $F_1,\dots,F_4$ for the four classes of independence
number $5$. All five graphs are \capeleven{} and $K_6$-free, and they partition
$E(K_{21})$, so
\begin{equation}\label{eq:MH-partition}
  e(H)+\sum_{i=1}^4 e(F_i)=\binom{21}{2}=210 .
\end{equation}
The proof shows $e(H)\ge 58$ and $e(F_i)\ge 38$, forcing equality throughout, and
then derives a contradiction from the rigid equality structure.

\subsection{Reduction to $21$ vertices}\label{ssec:MH-reduce}
Let $\chi$ be a balanced $5$-colouring of $K_{25}$ with a colour (say $0$) and a
$4$-set $T$ such that $\alpha(G_0-T)\le 4$. Put $S=V(K_{25})\setminus T$,
$\abs S=21$, and $H=G_0[S]$, $F_i=G_i[S]$ ($1\le i\le 4$). Then $\alpha(H)\le 4$
(an independent $5$-set of $H$ is one of $G_0-T$), and $\alpha(F_i)\le 5$
(Lemma~\ref{lem:elem}). All five are \capeleven{} and $K_6$-free
(Lemma~\ref{lem:elem}\ref{it:cap}, inherited by $S$). This is exactly the
configuration above, so Lemma~\ref{lem:MH} follows from its $21$-vertex form.

For each $v\in S$ set $W_v=S\setminus N_H[v]$. As in Section~\ref{sec:toolkit},
$\alpha(H[W_v])\le 3$ for every $v$ (an independent $4$-set in $W_v$ together with
$v$ is an independent $5$-set of $H$).

\subsection{Every ordinary colour has at least $38$ edges}\label{ssec:MH-Fi}
Fix an \emph{ordinary} colour $i\in\{1,2,3,4\}$.

\begin{claim}
$F_i$ is $K_5$-free.
\end{claim}
Indeed, a $K_5$ in $F_i$ is a $5$-set spanning only colour-$i$ edges, hence
spanning no $H$-edge, hence an independent $5$-set of $H$ --- contradicting
$\alpha(H)\le 4$.

Let $J_i=\overline{F_i}$ be the complement. Since $\alpha(F_i)\le 5$, $J_i$ is
$K_6$-free; since $\omega(F_i)\le 4$ ($F_i$ is $K_5$-free),
$\alpha(J_i)=\omega(F_i)\le 4$, so $J_i$ is not $5$-partite (five independent parts
would cover at most $5\cdot 4=20<21$ vertices). Brouwer's theorem
(Theorem~\ref{thm:brouwer}) with $r=5$, $n=21$ and the value~\eqref{eq:brouwer-values}
gives $e(J_i)\le 173$, hence
\begin{equation}\label{eq:MH-37}
  e(F_i)\ge\binom{21}{2}-173=210-173=37 .
\end{equation}

We exclude equality. Suppose $e(F_i)=37$. Then $J_i$ is an extremal graph in
Theorem~\ref{thm:brouwer} at $(r,n)=(5,21)$, so by the Kang--Pikhurko
classification (Corollary~\ref{cor:equality21}) $F_i$ contains disjoint sets
$A,B$ with $\abs A=5$, $\abs B=4$,
\[
  F_i[A]=K_5-xy\ \ (\text{some }x,y\in A),\qquad F_i[B]=K_4,\qquad e_{F_i}(A,B)=4 .
\]
(In Kang--Pikhurko's construction the part sizes of the extremal graph on the $20$
``old'' vertices are one of $(4,4,4,4,4)$, $(3,4,4,4,5)$, $(3,3,4,5,5)$; the latter
two contain a part of size $5$, which would be a $K_5$ in $F_i$, so only
$(4,4,4,4,4)$ survives, and translating that construction into $F_i=\overline{J_i}$
gives exactly the stated $A,B$ with $4$ cross-edges, independently of the internal
parameter of the construction.)

Now count. The non-edge $xy$ of $F_i[A]$ carries some colour $\ne i$; it cannot be
an $F_j$-edge for an ordinary $j\ne i$ with $A$ otherwise avoiding $H$, because $A$
must contain an $H$-edge (else $A$ is an independent $5$-set of $H$), and the only
non-$i$ edge inside $A$ is $xy$, so $xy$ is an $H$-edge. Consider the $20$ edges
between $A$ and $B$:
\begin{itemize}[nosep]
\item colour $i$ contributes exactly $4$ (given);
\item colour $0$ ($H$) contributes at least $5$: $B$ is $H$-independent (it is a
$K_4$ of $F_i$), so each $a\in A$ has an $H$-neighbour in $B$ (else $B\cup\{a\}$ is
an $H$-independent $5$-set);
\item each ordinary $j\ne i$ contributes at least $4$: for every pair $a,a'\in A$
the $6$-set $B\cup\{a,a'\}$ must contain a colour-$j$ edge, and neither $B$
($=K_4$ in $F_i$) nor the pair $aa'$ can supply it, so it runs from $\{a,a'\}$ to
$B$; hence at most one vertex of $A$ has no colour-$j$ neighbour in $B$, giving
$e_j(A,B)\ge 4$.
\end{itemize}
Edges of distinct colours are distinct, so the $A$--$B$ edges number at least
$4+5+3\cdot 4=21>20$, a contradiction. Therefore
\begin{equation}\label{eq:MH-38}
  e(F_i)\ge 38\qquad(1\le i\le 4).
\end{equation}

\subsection{If $H$ is $K_5$-free then $e(H)\ge 58$}\label{ssec:MH-58}
Assume for now that $H$ is $K_5$-free, so $\omega(H)\le 4$. Each $W_v$ has
$\alpha\le 3$, $\omega\le 4$ and is \capeleven, so $e(H[W_v])\ge L(20-d_v)$
(Lemma~\ref{lem:Ltable}), finite only for $20-d_v\le 20$. Feeding
$e(H[W_v])\ge L(20-d_v)$ and Lemma~\ref{lem:nbhd} into~\eqref{eq:41}--\eqref{eq:42}
(with $q=21$) yields $\sum_v\Psi(d_v)\le 0$, where
\[
  \Psi(d)=L(20-d)+d^2-\tfrac{21}{2}d-u(d).
\]
Direct substitution gives $\Psi(d)\ge 52-\tfrac{19}{2}d$ for $0\le d\le 20$; the
successive slacks $\Psi(d)-(52-\tfrac{19}{2}d)$ are
\[
  32,\,21,\,11,\,4,\,1,\,0,\,0,\,2,\,6,\,14,\,23,\,34,\,48,\,63,\,79,\,97,\,117,\,139,\,163,\,188,\,214,
\]
all nonnegative. Hence
\[
  0\ge\sum_v\Psi(d_v)\ge 21\cdot 52-\tfrac{19}{2}\sum_v d_v=1092-19\,e(H),
\]
so $e(H)\ge\lceil 1092/19\rceil=58$:
\begin{equation}\label{eq:MH-58}
  H\text{ $K_5$-free}\ \Longrightarrow\ e(H)\ge 58 .
\end{equation}

\subsection{$H$ is $K_5$-free}\label{ssec:MH-K5free}
Suppose, for contradiction, that $H$ contains a $K_5$, say on the vertex set $Q$,
and put $X=S\setminus Q$, $\abs X=16$.

By \capeleven{} (Lemma~\ref{lem:K5indeg}), every $x\in X$ has at most one
$H$-neighbour in $Q$. If $X$ contained an $H$-independent $4$-set, its four
vertices would have at most $4$ neighbours in $Q$, so some $q\in Q$ would be
$H$-nonadjacent to all of them, giving an independent $5$-set of $H$; thus
\begin{equation}\label{eq:MH-X-alpha}
  \alpha(H[X])\le 3 .
\end{equation}

We claim $H[X]$ is $K_5$-free. If not, let $R\subseteq X$ be an $H$-$K_5$ and
$L=X\setminus R$, $\abs L=11$. An independent triple in $L$, together with a vertex
of $Q$ and a vertex of $R$ chosen non-adjacent to it (possible since \capeleven{}
makes the $Q$--$R$ and triple--$Q$, triple--$R$ adjacencies sparse), would form an
$H$-independent $5$-set; hence $\alpha(H[L])\le 2$. But $L$ is a \capeleven{} graph
on $11$ vertices with $\alpha\le 2$, and no such graph exists by
Fact~\ref{sat:nonex11} (Remark~\ref{rem:nonex-mono}). So $H[X]$ contains no second
$K_5$, i.e.\ $H[X]$ is $K_5$-free.

\begin{remark}
This is the one place where the \emph{$\omega$-free} form of the nonexistence fact
is essential. The original draft bounded $\omega(H[L])\le 4$ and invoked a value
$M(11)=35$; but $M(11)$ is in fact $+\infty$ (no such graph exists at all,
Fact~\ref{sat:nonex11}), and the clean statement is that $L$ cannot exist for any
clique number. This makes the second-$K_5$ case vacuous outright.
\end{remark}

Since $H[X]$ is $K_5$-free with $\alpha\le 3$ and is \capeleven, Lemma~\ref{lem:Ltable}
gives
\begin{equation}\label{eq:MH-X-46}
  e(H[X])\ge L(16)=46 .
\end{equation}
Now bound the ordinary colours on $X$. For each $i$, $F_i[X]$ is $K_4$-free (a $K_4$
of $F_i$ inside $X$ would be an $H[X]$-independent $4$-set, contradicting
\eqref{eq:MH-X-alpha}). Its complement is $K_6$-free (as $\alpha(F_i[X])\le 5$) and
not $5$-partite ($\alpha=\omega(F_i[X])\le 3$, and $5\cdot 3=15<16$), so Brouwer's
theorem with $(r,n)=(5,16)$ and~\eqref{eq:brouwer-values} gives
$e(\overline{F_i[X]})\le 100$, hence
\begin{equation}\label{eq:MH-X-20}
  e(F_i[X])\ge\binom{16}{2}-100=120-100=20 .
\end{equation}
Summing~\eqref{eq:MH-X-46} and~\eqref{eq:MH-X-20} over the five colours inside $X$,
\[
  \binom{16}{2}=120\;\ge\; e(H[X])+\sum_{i=1}^4 e(F_i[X])\;\ge\; 46+4\cdot 20=126,
\]
a contradiction. Therefore $H$ is $K_5$-free.

\subsection{All bounds are equalities}\label{ssec:MH-eq}
By~\eqref{eq:MH-38}, \eqref{eq:MH-58} (now unconditional, as $H$ is $K_5$-free) and
the edge partition~\eqref{eq:MH-partition},
\[
  210=e(H)+\sum_{i=1}^4 e(F_i)\ge 58+4\cdot 38=210 .
\]
Hence
\begin{equation}\label{eq:MH-equalities}
  e(H)=58,\qquad e(F_i)=38\quad(1\le i\le 4).
\end{equation}

\subsection{Excluding the equality case}\label{ssec:MH-endgame}
We now derive a contradiction from~\eqref{eq:MH-equalities}.

\paragraph{Minimum degree.}
Suppose $d_H(v)=4$ for some $v$. Then $\abs{W_v}=16$, $\alpha(H[W_v])\le 3$,
$\omega(H[W_v])\le 4$, so $e(H[W_v])\ge L(16)=46$; and each $F_i[W_v]$ is $K_4$-free
(an $F_i$-$K_4$ in $W_v$ with $v$ is an $H$-independent $5$-set), giving
$e(F_i[W_v])\ge 20$ as in~\eqref{eq:MH-X-20}. Then the $\binom{16}2=120$ edges of
$K_{16}$ on $W_v$ would satisfy $120\ge 46+4\cdot 20=126$, impossible. The same count
(with the larger floors $L(17),\dots,L(19)\ge 46$) excludes $d_H(v)\le 3$ a
fortiori. Hence
\begin{equation}\label{eq:MH-delta5}
  \delta(H)\ge 5 .
\end{equation}
Since $2e(H)=116<6\cdot 21=126$, some vertex has degree exactly $5$. Fix such a
vertex $v$ and put
\[
  A=N_H(v),\quad Q=A\cup\{v\},\quad W=S\setminus Q,\qquad \abs A=5,\ \abs Q=6,\ \abs W=15 .
\]
Write $\rho=e_H(A)$,\footnote{The source documents in \texttt{review\_queue/} name
this quantity $r$; we rename it to $\rho$ to avoid shadowing the global colour count
$r=5$ (and correspondingly write $\mu_x$ below for the multiplicity they call
$\rho_x$).} $w=e_H(W)$, $c=e_H(A,W)$. Since $\alpha(H[W])\le 3$,
$\omega(H[W])\le 4$ and $W$ is \capeleven, $w\ge L(15)=38$; and \capeleven{} on $Q$
gives $\rho\le 6$.

\paragraph{The $(\rho,w,c)$ system.}
Let $a=\sum_{x\in A}(d_H(x)-5)\ge 0$ (by~\eqref{eq:MH-delta5}) and $\eta=w-38\ge 0$.
Summing $H$-degrees over $A$ counts each edge $xv$ once, each internal $A$-edge
twice, and each $A$--$W$ edge once (a vertex of $A$ has no other neighbours, as
$V(H)=Q\cup W$ and $Q=A\cup\{v\}$), so
$\sum_{x\in A}d_H(x)=\abs A+2\rho+c=5+2\rho+c$. On the other hand
$\sum_{x\in A}d_H(x)=25+a$ by the definition of $a$. Hence
\[
  25+a=5+2\rho+c,\qquad\text{so}\qquad c=20+a-2\rho .
\]
Using $e(H)=58=e_H(Q)+w+e_H(Q,W)=(\rho+5)+w+c$ (as $v$ has no neighbour in $W$, so
$e_H(Q,W)=e_H(A,W)=c$) and $w=38+\eta$,
\[
  58=(\rho+5)+(38+\eta)+(20+a-2\rho)=63-\rho+a+\eta,\qquad\text{i.e.}\qquad
  a+(6-\rho)+\eta=1 .
\]
Since $a,\eta\ge 0$ and $\rho\le 6$, the only solutions are
\begin{equation}\label{eq:MH-rwc}
  (\rho,w,c)\in\{(5,38,10),\ (6,38,9),\ (6,39,8)\}.
\end{equation}

\paragraph{A low ordinary colour on $W$, and $c\ge 10$.}
The four ordinary colours partition the non-$H$ edges of $K_{15}$ on $W$, so
$\sum_{i=1}^4 e(F_i[W])=\binom{15}{2}-w=105-w\in\{66,67\}$. Some ordinary colour,
say $i$, has $e(F_i[W])\le 16$. Also $F_i[W]$ is $K_4$-free (as $\alpha(H[W])\le 3$),
so its complement is $K_6$-free and, by the complement-Tur\'an bound for $\alpha\le
3$, $e(F_i[W])\ge L(15)$'s ordinary floor $=15$. Two cases:
\begin{itemize}[nosep]
\item If $e(F_i[W])=15$, then $\overline{F_i[W]}$ attains the Tur\'an bound
$t_3(15)=90$, and by uniqueness of the extremal graph $F_i[W]=5K_3$.
\item If $e(F_i[W])=16$, then $e(\overline{F_i[W]})=\binom{15}2-16=89$; a $K_6$-free
non-$5$-partite graph on $15$ vertices has at most $t_5(15)-3+1=88<89$ edges
(Theorem~\ref{thm:brouwer}), so $\overline{F_i[W]}$ is $5$-partite, and since
$\omega(F_i[W])\le 3$ all five parts have size $3$.
\end{itemize}
In either case $W=C_1\dot\cup\cdots\dot\cup C_5$ where each $C_j$ is a colour-$i$
triangle, hence $H$-independent. For each $j$ set $X_j=\{q\in Q:e_H(q,C_j)=0\}$. If
two vertices of $X_j$ were $H$-nonadjacent, they would form, with $C_j$, an
$H$-independent $5$-set; so $X_j$ is an $H$-clique, and since $H$ is $K_5$-free,
$\abs{X_j}\le 4$. Thus at least two vertices of $Q$ send an $H$-edge into each
$C_j$, whence
\begin{equation}\label{eq:MH-c10}
  c=e_H(Q,W)\ge 2\cdot 5=10 .
\end{equation}
This excludes the last two cases of~\eqref{eq:MH-rwc}, leaving $(\rho,w,c)=(5,38,10)$.

\paragraph{The triangle-cover contradiction.}
Equality in~\eqref{eq:MH-c10} means that for each $j$: exactly two vertices of $Q$
send an $H$-edge into $C_j$, each exactly one such edge, and the other four vertices
form the $H$-$K_4$ $X_j$. Since $v$ has no $H$-neighbour in $W$, $v\in X_j$ for every
$j$. Let $F=H[A]$; as $\rho=5$, $e(F)=5$. Write $X_j=\{v\}\cup T_j$, so $T_j$ is a
triangle of $F$. For $x\in A$ put $\mu_x=\abs{\{j:x\in T_j\}}$. The $H$-degree of
$x$ decomposes as
\[
  d_H(x)=\underbrace{1}_{xv}+\;d_F(x)\;+\underbrace{(5-\mu_x)}_{\text{edges into }W},
\]
because $x$ sends one $H$-edge into each $C_j$ with $x\notin T_j$ and none into the
$C_j$ with $x\in T_j$. By~\eqref{eq:MH-delta5}, $d_H(x)\ge 5$, so $\mu_x\le
d_F(x)+1$. But
\[
  \sum_{x\in A}\mu_x=5\cdot 3=15=2e(F)+5=\sum_{x\in A}(d_F(x)+1),
\]
so $\mu_x=d_F(x)+1\ge 1$ for every $x$; in particular the five triangles
$T_1,\dots,T_5$ cover all of $A$. This is impossible in a $5$-edge graph: any set of
triangles covering all $5$ vertices of $F$ uses at least $6$ edges (one triangle
uses $3$ edges and covers $3$ vertices; covering the remaining two needs either a
second triangle sharing an edge, which covers only one new vertex, or two triangles
adding at least $3$ more edges). This contradicts $e(F)=5$ and completes the proof
of Lemma~\ref{lem:MH}.\qed

\subsection{Why the argument does not apply at $n=24$}\label{ssec:MH-n24}
The proof uses $\abs S=21$ quantitatively; it must, because $4$-hitters do exist in
balanced colourings of $K_{24}$ (Appendix~\ref{app:phase}). Deleting a $4$-hitter
from a balanced $K_{24}$ leaves $20$ vertices, where the special colour can be
$4K_5$ ($40$ edges) and an ordinary $K_5$-free colour can be $5K_4$ ($30$ edges),
so the analogue of the equality count reads $40+4\cdot 30=160<\binom{20}2=190$ ---
no contradiction. Sharper: if the special colour contains a $K_5$, deleting it
leaves $15$ vertices with floors $30+4\cdot 15=90\le\binom{15}2=105$, again slack;
whereas at $21$ vertices the same deletion leaves $16$ and gives the decisive
$46+4\cdot 20=126>120=\binom{16}2$. The tension is exactly the odd part
$21=4\cdot 5+1$: the ``$+1$'' of the conjecture propagates down to force one
oversized part where $20=4\cdot 5$ splits evenly.

%% ===========================================================================
\section{Proof of Lemma \textsf{[MM]}}\label{sec:MM}
%% ===========================================================================

Suppose, for contradiction, that $G$ on $25$ vertices and a $5$-set $T$ satisfy the
hypotheses of Lemma~\ref{lem:MM}: $\alpha(G)\le 5$, $G$ is \capeleven,
$e(G)\le 60$, $\alpha(G-T)\le 4$, and $e(G[T])\le 6$. Put
\[
  H=G-T,\quad m=e(H),\quad x=e_G(T,H),\quad \sigma=e(G[T]).
\]
Then $\abs{V(H)}=20$, $\alpha(H)\le 4$, $\sigma\le 6$, and
\begin{equation}\label{eq:MM-budget0}
  m+x+\sigma\le 60 .
\end{equation}
Let $F$ denote the graph of \emph{non-edges} of $G[T]$ (a graph on the $5$ vertices
of $T$), and $\nu=e(F)=10-\sigma\ge 4$. Since $\alpha(G)\le 5$, $G$ has no
independent $6$-set --- a fact used constantly, together with
Lemma~\ref{lem:K5indeg} (\capeleven{} forces at most one edge from any vertex into a
$K_5$, and a matching between two disjoint $K_5$'s).

For a graph $Y$ with $\alpha(Y)\le 3$, $\omega(Y)\le 4$ and \capeleven{} we write
$L(\abs{V(Y)})$ for the edge floor of Lemma~\ref{lem:Ltable}; recall
$L(15)=38$, $L(16)=46$, $L(17)=53$, $L(18)=62$, and for $\le 12$ vertices the
complement-Tur\'an value.

\subsection{Peeling and the case split}
\begin{lemma}[Peeling]\label{lem:peel}
Let $Q_1,\dots,Q_k$ be disjoint $K_5$'s in $H$ and $R=H-\bigcup_i Q_i$. Then
$\alpha(R)\le 4-k$.
\end{lemma}
\begin{proof}
If $S\subseteq R$ were independent with $\abs S=5-k$, choose $q_i\in Q_i$
successively keeping all chosen vertices independent: at step $i$, the $\le 5-k$
vertices of $S$ and the $\le i-1$ earlier choices each rule out at most one vertex
of $Q_i$ (Lemma~\ref{lem:K5indeg}), so at most $(5-k)+(i-1)\le 4$ vertices are
ruled out and some $q_i\in Q_i$ remains. The result is an independent $5$-set of
$H$, contradicting $\alpha(H)\le 4$.
\end{proof}

Let $k$ be the maximum number of pairwise disjoint $K_5$'s in $H$. By
Lemma~\ref{lem:peel}, $k\le 4$; and $k=3$ is impossible, since then $R$ has
$\abs{V(R)}=5$ and $\alpha(R)\le 1$, i.e.\ $R$ is a fourth $K_5$. Hence
\[
  k\in\{0,1,2,4\},
\]
and we eliminate each case.

\subsection{The case $k=0$: $H$ is $K_5$-free}\label{ssec:MM-free}
Then every $W_v=H\setminus N_H[v]$ ($\abs{W_v}=19-d_v$) has $\alpha\le 3$,
$\omega\le 4$, and is \capeleven, so $e(H[W_v])\ge L(19-d_v)$. Writing $f(d)=b(d)$
for the neighbourhood bound of Lemma~\ref{lem:nbhd}, the
identities~\eqref{eq:41}--\eqref{eq:42} give $\sum_v g(d_v)\le 0$ with
\[
  g(d)=L(19-d)+d^2-f(d)-10d .
\]
Substituting $L$ and $f$, the quantity $2g(d)+17d-84$ equals, for $d=0,\dots,19$,
\[
  62,39,22,11,0,1,0,5,18,33,50,73,100,129,162,197,236,281,328,377,
\]
all nonnegative, so $g(d)\ge 42-\tfrac{17}2 d$. Then
$0\ge\sum_v g(d_v)\ge 20\cdot 42-\tfrac{17}2\cdot 2m=840-17m$, whence
\begin{equation}\label{eq:MM-m50}
  m\ge 50 .
\end{equation}

Now fix a non-edge $tt'$ of $G[T]$ and put $U_{tt'}=N_H(t)\cup N_H(t')$. An
independent $4$-set of $H-U_{tt'}$ together with $t,t'$ would be an independent
$6$-set of $G$; hence
\begin{equation}\label{eq:MM-Ualpha}
  \alpha(H-U_{tt'})\le 3 .
\end{equation}
Summing over the $\nu$ non-edges of $G[T]$, and using $d_F(t)\le 4$,
\begin{equation}\label{eq:MM-Usum}
  \sum_{tt'}\abs{U_{tt'}}\le\sum_{tt'}\bigl(d_H(t)+d_H(t')\bigr)
  =\sum_{t\in T}d_H(t)\,d_F(t)\le 4x .
\end{equation}
Because $m\le 60<62=L(18)$, we always have $\abs{U_{tt'}}\ge 3$ (else $H-U_{tt'}$
contains an $18$-vertex induced subgraph with $\alpha\le 3$, $\omega\le 4$,
\capeleven, forcing $\ge L(18)=62>m$ edges). Also~\eqref{eq:MM-budget0} gives
$x\le 60-m-\sigma=50+\nu-m$. We split on $m$.
\begin{itemize}[nosep]
\item If $m\ge 53$: from $\abs{U_{tt'}}\ge 3$ and~\eqref{eq:MM-Usum},
$3\nu\le 4x\le 4(50+\nu-m)\le 4(\nu-3)$, so $\nu\ge 12>10$, impossible.
\item If $m\in\{51,52\}$: now $\abs{U_{tt'}}\ge 4$ (a $17$-vertex residual would need
$L(17)=53>m$), so $4\nu\le 4x$, i.e.\ $\nu\le x\le 50+\nu-m\le\nu-1$, impossible.
\item If $m=50$: first $\delta(H)\ge 3$ (if $d(v)\le 2$ then
$m\ge d(v)+L(19-d(v))\ge\min\{73,63,55\}=55>50$). If some $\abs{U_{tt'}}\le 4$,
enlarge it to a $4$-set $U$; then $H-U$ has $\ge L(16)=46$ edges and the number of
edges meeting $U$ is at least $\sum_{u\in U}d(u)-\binom 42\ge 12-6=6$, so
$m\ge 46+6=52>50$, impossible. Hence $\abs{U_{tt'}}\ge 5$ for every non-edge, and
$5\nu\le 4x\le 4(50+\nu-50)=4\nu$, so $\nu\le 0$, impossible.
\end{itemize}
Thus $H$ is not $K_5$-free.

\subsection{The case $k=4$: four disjoint $K_5$'s}
Write $H=Q_1\cup Q_2\cup Q_3\cup Q_4$. By Lemma~\ref{lem:K5indeg} every cross-graph
$H[Q_i,Q_j]$ is a matching and every $t\in T$ has at most one neighbour in each
$Q_i$. Choose a non-edge $tt'$ of $G[T]$ (it exists as $\sigma\le 6<10$), and set
$A_i=Q_i\setminus(N(t)\cup N(t'))$, so $\abs{A_i}\ge 3$. There is no independent
transversal of $A_1,\dots,A_4$: such a transversal with $t,t'$ would be an
independent $6$-set.

\begin{claim}
If four parts each have size $\ge 3$, every cross-graph is a matching, and there is
no independent transversal, then every vertex of every part has a neighbour in each
other part.
\end{claim}
\begin{proof}
Suppose $v\in A_i$ has no neighbour in $A_j$. Delete from each of the two other
parts $A_k$ ($k\ne i,j$) the at most one neighbour of $v$; at least two vertices
remain in each. The matching between these two residual sets cannot cover all pairs,
so pick non-adjacent $y,z$, one in each, both non-adjacent to $v$. Then $y,z$
forbid at most two vertices of $A_j$ and $v$ forbids none, so
($\abs{A_j}\ge 3$) a fourth vertex of $A_j$ completes an independent transversal.
\end{proof}

Hence each cross-graph $H[A_i,A_j]$ saturates both sides, so has $\ge 3$ edges, and
\[
  e(H)\ge 4\binom 52+6\cdot 3=40+18=58 .
\]
By~\eqref{eq:MM-budget0}, $\sigma+x\le 2$. Now $G[T]$ has at most $2$ edges on $5$
vertices, so $T$ contains an independent triple $P$. Deleting from $Q_1,Q_2,Q_3$
all neighbours of $P$ removes at most $x\le 2$ vertices in total, leaving $\ge 3$ in
each clique; as the cross-graphs are matchings, a greedy independent transversal
$q_1\in Q_1$, $q_2\in Q_2$, $q_3\in Q_3$ exists, and $P\cup\{q_1,q_2,q_3\}$ is an
independent $6$-set of $G$. Contradiction.

\subsection{The case $k=1$: exactly one $K_5$}
Let $Q$ be the unique $K_5$ and $R=H-Q$, so $\abs{V(R)}=15$, $\alpha(R)\le 3$
(Lemma~\ref{lem:peel}), and $\omega(R)\le 4$ ($R$ is $K_5$-free, since a $K_5$ in
$R$ would be disjoint from $Q$). Hence
\begin{equation}\label{eq:MM-eR}
  e(R)=38+b,\qquad b\ge 0 .
\end{equation}
Set $p=e(Q,R)$; by Lemma~\ref{lem:K5indeg} each $r\in R$ has at most one neighbour
in $Q$. For $t_i\in T$, let $q_i$ be its (at most one) neighbour in $Q$, let
$Z_i=N_R(t_i)$, $d_i=\abs{Z_i}$, let $a$ be the number of $t_i$ with $q_i$ defined,
and $D=\sum_i d_i$. Counting edges,
\begin{equation}\label{eq:MM-budget1}
  b+p+a+D+\sigma\le 12 ,
\end{equation}
since $e(G)=e(Q)+e(R)+e(Q,R)+e(T,H)+e(G[T])=10+(38+b)+p+(a+D)+\sigma\le 60$.

For each edge $ij$ of $F$ (i.e.\ each non-edge $t_it_j$ of $G[T]$):
\begin{itemize}[nosep]
\item if $q_i,q_j$ are not two distinct vertices of $Q$, then
$\alpha(R-(Z_i\cup Z_j))\le 2$ (an independent triple there, with $t_i,t_j$, has at
most $3+1=4$ neighbours in $Q$, so misses a $Q$-vertex, giving an independent
$6$-set), and by Fact~\ref{sat:nonex11} this residual $\omega\le 4$, \capeleven{}
graph has $\le 10$ vertices; hence
\begin{equation}\label{eq:MM-14}
  d_i+d_j\ge \abs{Z_i\cup Z_j}\ge 5 ;
\end{equation}
\item if $q_i,q_j$ are distinct, let $P\subseteq R$ be the $p$ vertices with a
$Q$-neighbour; the same argument inside $R-(Z_i\cup Z_j\cup P)$ (the triple now has
no $Q$-neighbour) gives
\begin{equation}\label{eq:MM-15}
  d_i+d_j+p\ge 5 .
\end{equation}
\end{itemize}
Let $h$ be the number of $F$-edges joining two ``marked'' vertices (those with a
$Q$-neighbour) whose $Q$-neighbours are distinct. Summing~\eqref{eq:MM-14}
and~\eqref{eq:MM-15} over $E(F)$,
\[
  \sum_{ij\in E(F)}(d_i+d_j)\ge 5\nu-\min(p,5)\,h,
\]
while $\sum_{ij\in E(F)}(d_i+d_j)=\sum_i d_i\,d_F(i)\le 4D$, and $h\le\min\{\nu,\binom a2\}$.
From~\eqref{eq:MM-budget1}, $D\le \nu+2-p-a$. Hence a necessary condition is
\begin{equation}\label{eq:MM-16}
  5\nu-\min(p,5)\min\{\nu,\tbinom a2\}\le 4(\nu+2-p-a).
\end{equation}
An exhaustive scan of~\eqref{eq:MM-16} over $a\in\{0,\dots,5\}$, $p\ge 0$,
$\nu\in\{4,\dots,10\}$ (verified by exact arithmetic; see
Section~\ref{sec:certificates}) leaves exactly two families:
\[
  a\le 1\text{ and } p+a\le 1;\qquad\text{or}\qquad a=5\text{ and }p\in\{5,6,7\};
\]
the borderline $(p,a,\nu)=(4,5,10)$ is separately excluded, since then every pair of
$T$ has weight-sum $\ge 1$, so at most one $d_i$ is $0$ and $D\ge 4$, contradicting
$D\le\nu+2-p-a=3$.

\paragraph{Sub-case $p+a\le 1$.}
Every $F$-edge satisfies $d_i+d_j\ge 5$~\eqref{eq:MM-14}. If $F$ is not a star it
contains two disjoint edges, so $D\ge 10$ and hence $\nu\ge 8$
(via~\eqref{eq:MM-budget1}); but then $F=K_5$ minus $\le 2$ edges contains a
$5$-cycle, and summing $d_i+d_j\ge 5$ around it gives $2D\ge 25$, so $D\ge 13>\nu+2$,
impossible. Thus $F=K_{1,4}$, $\nu=4$, with centre $t$. Summing the four star
inequalities, $3\,d_R(t)+D\ge 20$; since $D\le\nu+2-p-a\le 6$, we get $d_R(t)\ge 5$
and at least one leaf $t_\ell$ has $d_\ell=0$. For that leaf, \eqref{eq:MM-14} gives
$\alpha(R-N_R(t))\le 2$. Consider $X=R+t$ (so $\abs{V(X)}=16$): then $\alpha(X)\le 3$
(an independent $4$-set through $t$ would give an independent triple in
$R-N_R(t)$), $X$ is \capeleven, and $e(X)=e(R)+d_R(t)\le 44$
(from~\eqref{eq:MM-budget1}). If $X$ is $K_5$-free then $e(X)\ge L(16)=46>44$,
impossible; if $X$ contains a $K_5$ it runs through $t$ (as $R$ is $K_5$-free), and
deleting it leaves $11$ vertices with $\omega\le 4$ whose independence number is at
most $2$ (an independent triple would, by Lemma~\ref{lem:K5indeg}, miss a vertex of
the deleted $K_5$, giving an independent $4$-set of $X$), contradicting
Fact~\ref{sat:nonex11}.

\paragraph{Sub-case $p\in\{5,6,7\}$, $a=5$.}
Now~\eqref{eq:MM-budget1} gives $D+\sigma\le 7-p$, so $D\le 2$ and
$\nu=10-\sigma\ge 8$. Delete $\bigcup_i Z_i$ (at most $D\le 2$ vertices) from $R$;
at least $13$ vertices remain, and by Fact~\ref{sat:nonex11} they cannot have
$\alpha\le 2$, so they contain an independent triple $S$ disjoint from all $Z_i$.
For every $F$-edge $ij$, $S\cup\{t_i,t_j\}$ is an independent $5$-set; to avoid an
independent $6$-set, its neighbours must cover all of $Q$, forcing
$\abs{N_Q(S)}=3$ and $\{q_i,q_j\}=Q\setminus N_Q(S)$ --- the same fixed $2$-element
set for every $F$-edge. (If $\abs{N_Q(S)}<3$ then no $F$-edge can cover $Q$, an
immediate independent $6$-set.) Thus $F$ embeds in a bipartite graph on the
$\le 5$ vertices of $T$ split by $q_i\in\{$the two fixed $Q$-vertices$\}$, so
$\nu=e(F)\le\lfloor 25/4\rfloor=6<8\le\nu$, a contradiction.

\subsection{The case $k=2$: exactly two $K_5$'s}\label{ssec:MM-two}
Let $Q_1,Q_2$ be the two cliques and $B$ the remaining $10$ vertices. By
Lemma~\ref{lem:peel}, $\alpha(B)\le 2$; and $\omega(B)\le 4$ (a $K_5$ in $B$ would
give three disjoint $K_5$'s). Hence
\begin{equation}\label{eq:MM-eB}
  e(B)=25+b,\qquad b\ge 0
\end{equation}
by $M(10)=25$ (Fact~\ref{sat:M10}; here $\omega\le 4$ is available, as required by
Remark~\ref{rem:nonex-mono}). Set
\[
  c=e(Q_1,Q_2)+e(B,Q_1\cup Q_2),\quad y=e(T,Q_1\cup Q_2),\quad D=e(T,B),
\]
so that $e(H)=45+b+c$ and~\eqref{eq:MM-budget0} becomes
\begin{equation}\label{eq:MM-budget2}
  b+c+y+D\le \nu+5 .
\end{equation}
For $t_i\in T$, $Z_i=N_B(t_i)$, $d_i=\abs{Z_i}$. Call an $F$-edge $ij$ \emph{unhit}
if $B-(Z_i\cup Z_j)$ contains a non-edge $uv$.

\paragraph{No unhit edges.}
For an unhit edge, $\{t_i,t_j,u,v\}$ is an independent $4$-set. Each of its four
vertices has at most one neighbour in each $Q_\ell$, so each $Q_\ell$ has $\ge 1$
missed vertex; if some missed vertex of $Q_1$ were non-adjacent to some missed
vertex of $Q_2$, those two would extend the $4$-set to an independent $6$-set.
Since $Q_1$--$Q_2$ edges form a matching (Lemma~\ref{lem:K5indeg}), this forces
exactly one missed vertex in each $Q_\ell$, adjacent to each other; consequently
$t_i,t_j$ each have $2$ neighbours in $Q_1\cup Q_2$, the vertices $u,v$ send $4$
edges from $B$ into $Q_1\cup Q_2$, and one $Q_1$--$Q_2$ edge is used. If there are
$h$ unhit edges meeting $\rho$ vertices of $T$, then
\begin{equation}\label{eq:MM-22}
  y\ge 2\rho,\qquad c\ge 5 .
\end{equation}
Every \emph{hit} $F$-edge has $B-(Z_i\cup Z_j)$ a clique, hence of size $\le 4$
($\omega(B)\le 4$), so $\abs{Z_i\cup Z_j}\ge 6$; summing,
$6(\nu-h)\le\sum_{ij}(d_i+d_j)\le 4D$. Combined with~\eqref{eq:MM-budget2}
and~\eqref{eq:MM-22} (which give $D\le\nu-2\rho$),
\begin{equation}\label{eq:MM-24}
  \nu+4\rho\le 3h .
\end{equation}
Since $h\le\binom\rho2$, \eqref{eq:MM-24} is impossible for $\rho\in\{2,3,4\}$
(using $\nu\ge 4$). For $\rho=5$ it forces $\nu=10$, $D=0$, $c=5$, $y=10$; then $T$
is independent and has no edge to $B$, so any $b\in B$ makes $T\cup\{b\}$ an
independent $6$-set. Hence $h=0$: every $F$-edge is hit, and
\begin{equation}\label{eq:MM-25}
  d_i+d_j\ge \abs{Z_i\cup Z_j}\ge 6\qquad\text{for every }ij\in E(F).
\end{equation}

\paragraph{Classifying $\nu$.}
From~\eqref{eq:MM-budget2}, $D\le\nu+5$.
\begin{itemize}[nosep]
\item $\nu\in\{5,6\}$: $F$ has two disjoint edges, so $D\ge 12>\nu+5$, impossible.
\item $\nu\in\{8,9\}$: $F$ contains a $5$-cycle, so $2D\ge 30$, $D\ge 15>\nu+5$,
impossible.
\item $\nu=10$: summing~\eqref{eq:MM-25} over all $10$ pairs, $4D\ge 60$, so $D=15$
and all $d_i=3$; then $\abs{Z_i\cup Z_j}=6$ forces the five $3$-sets $Z_i$ pairwise
disjoint, impossible in a $10$-set.
\item $\nu=7$ \textup{(the repaired case)}: $F$ has two disjoint edges, so $D\ge 12$;
with $D\le\nu+5=12$, $D=12$. The weight vector $(d_1,\dots,d_5)$ is nonnegative,
sums to $12$, and has at least $7$ pairs of sum $\ge 6$ (one per $F$-edge); the only
such vector is $(0,0,0,6,6)$. So the two weight-$6$ vertices are $F$-adjacent to each
other and to the three weight-$0$ vertices. Pick a weight-$6$ vertex $t_i$
($\abs{Z_i}=6$) and a weight-$0$ vertex $t_j$ ($Z_j=\varnothing$); their $F$-edge is
hit, so $B-Z_i$ is a clique, of size $10-6=4$, i.e.\ a $K_4$. Then $X=B+t_i$ has
$\abs{V(X)}=11$, is \capeleven, and has $\alpha(X)\le 2$ (an independent triple in
$X$ would be a triple in $B$, impossible, or $\{t_i,u,v\}$ with $u,v\in B-Z_i=K_4$
adjacent, impossible). This contradicts Fact~\ref{sat:nonex11}.
\item $\nu=4$: $F=K_{1,4}$ (else two disjoint edges give $D\ge 12>9$), with centre
$t_0$. If a leaf has weight $0$, then $B-N_B(t_0)$ is a clique and $X=B+t_0$ is an
$11$-vertex \capeleven{} graph with $\alpha\le 2$ (as in the $\nu=7$ case),
contradicting Fact~\ref{sat:nonex11}. So all four leaves have positive weight; then
\eqref{eq:MM-25} and $D\le 9$ force
\[
  d_B(t_0)=5,\quad d_B(t_i)=1\ (1\le i\le 4),\quad D=9,
\]
and~\eqref{eq:MM-budget2} forces $b=c=y=0$, so $e(B)=25$ and there are no edges
between $T$ and $Q_1\cup Q_2$, none between $B$ and $Q_1\cup Q_2$, and none between
$Q_1$ and $Q_2$. We finish this last case below.
\end{itemize}

\paragraph{The $\nu=4$ endgame.}
Put $Z=N_B(t_0)$ and $C=B-Z$, so $\abs Z=\abs C=5$, and let $z_i$ be the unique
$B$-neighbour of leaf $t_i$. Since $Z\cup\{z_i\}$ meets every non-edge of $B$ (the
$F$-edge $t_0t_i$ is hit), $B-(Z\cup\{z_i\})$ is a clique; as $\omega(B)\le 4$ this
forces $z_i\in C$ and
\begin{equation}\label{eq:MM-30}
  B[C\setminus\{z_i\}]=K_4 .
\end{equation}
For every $v\in B$, its non-neighbours form a clique (as $\alpha(B)\le 2$) of size
$\le 4$, so $d_B(v)\ge 5$; since $e(B)=25$ has average degree $5$, $B$ is
$5$-regular. Comparing degree sums of $Z$ and $C$ (each $=25$ minus the common
cross-count) gives $e(B[Z])=e(B[C])$. The $6$-set $\{t_0\}\cup Z$ gives
$5+e(B[Z])\le 11$, so $e(B[Z])\le 6$; and~\eqref{eq:MM-30} gives $e(B[C])\ge 6$.
Hence
\[
  e(B[Z])=e(B[C])=6 .
\]
Thus $z_i$ is isolated in $B[C]$ (the six edges of $C\setminus\{z_i\}\cong K_4$
exhaust $e(B[C])$), so its five neighbours all lie in $Z$: $z_i$ is adjacent to all
of $Z$. Finally, for any $z\in Z$ the $6$-set $\{t_0,z_i\}\cup(Z\setminus\{z\})$
spans $4+4+(6-d_{B[Z]}(z))=14-d_{B[Z]}(z)$ edges, so \capeleven{} gives
$d_{B[Z]}(z)\ge 3$ for every $z\in Z$; then $12=2e(B[Z])=\sum_{z\in Z}d_{B[Z]}(z)\ge
15$, a contradiction.

\medskip
All four cases $k\in\{0,1,2,4\}$ are impossible, so no such $G$ exists. This proves
Lemma~\ref{lem:MM}.\qed

\begin{remark}
The two boundary hypotheses of Lemma~\ref{lem:MM} enter essentially. The bound
$e(G[T])\le 6$ gives $\nu=10-e(G[T])\ge 4$, supplying the non-edges of $T$ that
drive every weighted hitting argument; and $e(G)\le 60$ supplies the sharp budgets
\eqref{eq:MM-budget0}, \eqref{eq:MM-budget1}, \eqref{eq:MM-budget2}, forcing exact
equalities in the final $k=2$, $\nu=4$ endgame.
\end{remark}

%% ===========================================================================
\section{The affine near-miss never extends}\label{sec:affine}
%% ===========================================================================

We now prove Theorem~\ref{thm:affine}, which is independent of
Sections~\ref{sec:toolkit}--\ref{sec:MM} and holds for all prime powers at once. It
says that the only balanced colourings of $K_{r^2}$ known in the literature ---
the affine ones --- are a genuine dead end for a counterexample: none extends by a
single vertex.

Fix a prime power $q\ge 3$ and the affine plane $\mathrm{AG}(2,q)$: $q^2$ points,
lines of $q$ points, and $q+1$ \emph{parallel classes} (directions), each a
partition of the points into $q$ disjoint lines. We use three standard facts:
each parallel class partitions the points into $q$ lines of size $q$ (F1); two lines
from different classes meet in exactly one point (F2); every pair of points lies on a
unique line, and ``collinear in a fixed direction'' is an equivalence relation whose
classes are the lines of that direction (F3).

Form the colouring $\mathcal C$ of $K_{q^2}$: colour each edge by the direction of
the line through its two endpoints, then merge two chosen directions $d_1,d_2$ into
colour $0$, letting the remaining directions $e_1,\dots,e_{q-1}$ be colours
$1,\dots,q-1$. Each line is monochromatic (its $\binom q2$ edges all have its own
direction). That $\mathcal C$ is balanced is immediate: any $(q+1)$-set meets some
line of every direction twice (pigeonhole: $q+1$ points into $q$ parallel lines), so
sees every colour.

\begin{proof}[Proof of Theorem~\ref{thm:affine}]
Suppose $\varphi\colon V(K_{q^2})\to\{0,\dots,q-1\}$ gives the colours of the edges
from a new vertex $x$, and the extension is balanced. A $(q+1)$-set through $x$ is
$\{x\}\cup F$ with $\abs F=q$; balancedness means every colour missing inside $F$
appears among $\{\varphi(p):p\in F\}$.

\smallskip
\noindent\emph{Step 1: each colour class of $\varphi$ has exactly $q$ points.}
Every line $L$ is monochromatic in $\mathcal C$, so as a $q$-set it misses the
other $q-1$ colours; hence $\varphi$ must realise all $q-1$ of them on the $q$
points of $L$. Thus for any colour $c$ and any direction $D$ whose colour is not
$c$, colour $c$ appears on \emph{every} one of the $q$ disjoint lines of $D$ (F1),
so $\abs{\varphi^{-1}(c)}\ge q$. Every colour has such a direction: colour $0$ is
exempt only on $d_1,d_2$ (leaving $q-1\ge 2$ directions), and each colour $c\ne 0$
is exempt only on $e_c$ (leaving $q$). Since $\sum_c\abs{\varphi^{-1}(c)}=q^2$ over
$q$ colours, all inequalities are equalities: $\abs{\varphi^{-1}(c)}=q$ for every
$c$.

\smallskip
\noindent\emph{Step 2: for $c\in\{1,\dots,q-1\}$, $T_c:=\varphi^{-1}(c)$ is a line
of direction $e_c$.}
Colour $c$ is exempt only on lines of direction $e_c$. So for each of the other $q$
directions ($d_1,d_2$, and $e_j$ with $j\ne c$), every line has colour $\ne c$ and
hence meets $T_c$. As $\abs{T_c}=q$ and there are $q$ disjoint lines per direction,
$T_c$ meets each such line exactly once; thus no two points of $T_c$ are collinear
in any non-exempt direction. Every pair of points is collinear in exactly one of the
$q+1$ directions, and $q$ of them are excluded, so every pair of $T_c$ is collinear
in $e_c$. Being in a common $e_c$-line is an equivalence relation (F3), so all $q$
points of $T_c$ lie in one $e_c$-line; as lines have $q$ points, $T_c$ \emph{is}
that line.

\smallskip
\noindent\emph{Step 3: contradiction.}
Since $q\ge 3$ there are $q-1\ge 2$ unmerged colours; take distinct $c,c'$. By
Step~2, $T_c,T_{c'}$ are lines of the distinct directions $e_c\ne e_{c'}$, so by
(F2) they meet in a point $p$. Then $\varphi(p)=c$ and $\varphi(p)=c'$, impossible
as colour classes are disjoint.
\end{proof}

For $q=2$ the argument breaks (only one unmerged colour, so Step~3 is vacuous), and
indeed $\Ncap(2)=5$ is realised by the pentagon, not by an affine construction. For
$q=5$ the theorem is confirmed computationally: all $15$ ways of choosing the merged
pair give a one-vertex extension problem that a SAT solver reports unsatisfiable
(Section~\ref{sec:certificates}).

%% ===========================================================================
\section{Computational certificates}\label{sec:certificates}
%% ===========================================================================

The proof uses computation in two sharply separated ways. \emph{All arithmetic}
--- the Tur\'an values~\eqref{eq:brouwer-values}, the $M$- and $L$-tables
(Lemmas~\ref{lem:Mtable}--\ref{lem:Ltable}), the affine-bound and slack
computations, the $g$-table of Section~\ref{ssec:MM-free}, and the case
scan~\eqref{eq:MM-16} --- is elementary and finite, and is confirmed by exact
integer recomputation (no floating point, no solver trust); the scripts are
\texttt{tools/verify\_gpt\_arith.py} (the $L$-recursion and slacks) and
\texttt{tools/verify\_gpt\_tables.py} (small SAT cross-checks of the tables). Only
four genuinely combinatorial facts about small graphs --- Fact~\ref{fact:sat} ---
are established by SAT, and each comes with an independently checkable
unsatisfiability certificate.

\subsection{The four primitives, as CNF}
Encode a graph on $s$ vertices by one Boolean variable $x_{\{a,b\}}$ per pair
$a<b$, true iff $\{a,b\}$ is an edge. The three graph properties become clause
families:
\begin{itemize}[nosep]
\item $\alpha\le 2$: every $3$-set spans an edge --- one clause
$\bigvee x_{\{a,b\}}$ per triple;
\item $\omega\le 4$: every $5$-set omits an edge --- one clause of $10$ negated
literals per $5$-set;
\item \capeleven: every $6$-set spans $\le 11$ of its $15$ pairs --- for each
$6$-set, one clause $\bigvee_{p\in P}\lnot x_p$ over each $12$-subset $P$ of its
pairs.
\end{itemize}
For the ``$\ge k$ edges'' facts we add a sequential-counter (Sinz) encoding of
``at most $k-1$ edges'' and ask for unsatisfiability. The four instances are then:
\[
  \renewcommand{\arraystretch}{1.2}
  \begin{array}{lllll}
    \textbf{name} & \textbf{statement} & \textbf{vars} & \textbf{clauses} & \textbf{verdict}\\\hline
    \textsf{nonex11} & \text{Fact~\ref{sat:nonex11}: no such graph, }s{=}11 & 55 & 210{,}375 & \text{UNSAT}\\
    \textsf{nonex12} & \text{Fact~\ref{sat:nonex12}: no such graph, }s{=}12 & 66 & 420{,}640 & \text{UNSAT}\\
    \textsf{M9$\ge$19} & \text{Fact~\ref{sat:M9}: no such graph with }\le 18\text{ edges} & 360 & 39{,}078 & \text{UNSAT}\\
    \textsf{M10$\ge$25} & \text{Fact~\ref{sat:M10}: no such graph with }\le 24\text{ edges} & 549 & 96{,}927 & \text{UNSAT}
  \end{array}
\]
(For \textsf{nonex11}/\textsf{nonex12} no $\omega$-clauses are present --- the facts
are clique-number-free, per Remark~\ref{rem:nonex-mono}. The variable counts $55$
and $66$ are $\binom{11}2$ and $\binom{12}2$; the extra variables of the $M$
instances are Sinz counters.) The clause counts are reproducible: e.g.\ \textsf{nonex11}
has $\binom{11}3=165$ triple-clauses and $\binom{11}6\binom{15}{12}=210{,}210$
cap-clauses, totalling $210{,}375$.

\subsection{Two independent verification paths}
Each instance is refuted with a resolution-style proof that is checked by a program
far simpler than any SAT solver.

\emph{(i) DRAT, checked by \texttt{drat-trim}.} A CDCL solver (kissat) emits a binary
DRAT proof (\texttt{data/sat/prim\_*.drat}); the trusted checker \texttt{drat-trim}
(built in \texttt{ext/drat-trim}) replays it and reports \verb|s VERIFIED| for all
four (\texttt{data/sat/prim\_*.check}). For example \textsf{nonex11} verifies
$158{,}868$ of $210{,}375$ clauses and $1{,}573{,}853$ lemmas in the unsatisfiability
core via $22{,}469{,}404$ resolution steps, with $0$ RAT lemmas.

\emph{(ii) LRAT, checked inside Lean.} For the formalisation (Section~\ref{sec:lean})
each instance is re-solved with cadical (with inprocessing disabled, so no fresh
variables are introduced), its LRAT proof is trimmed to consecutive clause
identifiers, and Lean's kernel-verified certificate checker
(\texttt{Std.Tactic.BVDecide}'s \texttt{verifyCert}) confirms the refutation by
reflection. This makes the SAT facts part of the Lean proof rather than an external
input, at the cost of one reflection axiom (\texttt{ofReduceBool} via
\texttt{native\_decide}; see Section~\ref{sec:lean}).

The two paths use different solvers, different proof formats, and different checkers,
and agree.

\subsection{Regenerating the certificates}
The DRAT certificates regenerate from the CNFs with
\begin{quote}\ttfamily
kissat --unsat data/sat/prim\_NAME.cnf data/sat/prim\_NAME.drat\\
ext/drat-trim/drat-trim data/sat/prim\_NAME.cnf data/sat/prim\_NAME.drat
\end{quote}
which prints \verb|s VERIFIED|; the CNFs themselves regenerate from the encoder in
\texttt{tools/}. The Lean LRAT certificates regenerate via the pipeline documented
in \texttt{FORMAL.md} (item F3): \texttt{cadical --lrat --no-binary
--inprocessing=false}, then the trimming step, then \texttt{lake build}.

\subsection{Corroborating computations (not part of the proof)}
Two further computations support, but are not required by, the results. First, the
referee \texttt{tools/verify.py} confirms that the $15$ affine colourings of
$K_{25}$ (all merge choices) are balanced --- ``all $\binom{25}{6}=177{,}100$
$6$-subsets see all $5$ colours'' --- witnessing $\Ncap(5)\ge 25$. Second,
\texttt{tools/extend.py} confirms that each of those colourings, and every one of
$120+$ sampled non-affine balanced colourings of $K_{25}$, has an unsatisfiable
one-vertex extension problem --- the computational shadow of Theorem~\ref{thm:affine}
and Theorem~\ref{thm:main} (Appendix~\ref{app:phase}).

%% ===========================================================================
\section{Lean formalisation}\label{sec:lean}
%% ===========================================================================

The entire argument of Sections~\ref{sec:reduction}--\ref{sec:MM} has been
formalised in Lean~4 with Mathlib \texttt{v4.30.0} (toolchain
\texttt{leanprover/lean4:v4.30.0}), in the project \texttt{lean617/} of the
accompanying repository. We describe what is proved --- the headline theorem now
carries no mathematical hypothesis --- and its axiom profile.

\subsection{What is formalised}
The colouring is modelled as \texttt{c : Sym2 (Fin 26) $\to$ Fin 5}, matching the
upstream statement (below). The following are all \texttt{sorry}-free:
\begin{itemize}[nosep]
\item the reduction (Lemma~\ref{lem:chain}), as \texttt{chain\_deduction : MH2 $\to$
MM $\to$ Main};
\item Lemma~\ref{lem:MH} (\textsf{[MH$''$]}), as \texttt{lemma\_MH2\_of}, including the
repaired $\omega$-free argument of Section~\ref{ssec:MH-K5free} and the endgame of
Section~\ref{ssec:MH-endgame};
\item Lemma~\ref{lem:MM} (\textsf{[MM]}), as \texttt{lemma\_MM\_of}, including the
repaired $\nu=7$ case of Section~\ref{ssec:MM-two};
\item the counting identities~\eqref{eq:41}--\eqref{eq:42}, the neighbourhood
bound~(Lemma~\ref{lem:nbhd}), and the $L$-table recursion~(Lemma~\ref{lem:Ltable});
\item the four SAT primitives (Fact~\ref{fact:sat}), discharged as
\texttt{primFacts : PrimFacts} by reflecting the cadical LRAT certificates through
Lean's verified certificate checker (Section~\ref{sec:certificates}, path~(ii));
\item Brouwer's edge bound itself (the ``saving'' half of Theorem~\ref{thm:brouwer}),
as \texttt{kp\_saving} --- the full Kang--Pikhurko induction (Case~A, the good/bad
dichotomy, Lemma~3, and the singleton guard) is formalised and axiom-clean, so the
bound is now \emph{proved}, not assumed;
\item the Kang--Pikhurko \emph{equality} classification at $(5,21)$
(Corollary~\ref{cor:equality21}), as
\texttt{kp\_equality\_\allowbreak classification\_\allowbreak proven}, via
the cone-descent discharge of Section~\ref{ssec:lean-brouwer} --- so this last classical
input is now \emph{proved} as well, and the final theorem needs no mathematical
hypothesis.
\end{itemize}
The headline theorem and its upstream-shaped corollary are
\begin{quote}\ttfamily\small
theorem erdos\_617\_r5\_unconditional : Main\\
theorem erdos\_617\_r5\_upstream\_unconditional\\
\phantom{xx}\{V\} [Fintype V] [DecidableEq V] (hV : Fintype.card V = 5\textasciicircum2 + 1)\\
\phantom{xx}(coloring : Sym2 V $\to$ Fin 5) : \ldots
\end{quote}
where \texttt{Main} is the $\mathrm{Fin}\,26$ statement of Theorem~\ref{thm:main}
and the \texttt{\_upstream\_} form is its transport to the shape of
\texttt{erdos\_617} at $r=5$ in the google-deepmind \texttt{formal-conjectures}
repository, over an arbitrary $26$-element vertex type~\cite{FC}. Neither carries a
mathematical hypothesis: each is obtained from the conditional
\texttt{erdos\_617\_r5 (h : KPEqualityClassification)} by supplying the now-proved
\texttt{kp\_equality\_classification\_proven} (Section~\ref{ssec:lean-brouwer}), which
in turn feeds \texttt{brouwerFacts\_of} to assemble the full \texttt{BrouwerFacts}
interface from the proved \texttt{kp\_saving} and the proved classification. The
conditional \texttt{erdos\_617\_r5} and \texttt{erdos\_617\_r5\_upstream} remain
exported, for a reader who prefers to audit the classification separately and supply it
as a hypothesis. The project \texttt{lake build} is clean, and
\texttt{erdos\_617\_r5\_unconditional} together with everything it imports contains no
admitted proof: \texttt{\#print axioms} lists no \texttt{sorryAx} (see the axiom profile
below; the definitive ``no admitted proof'' check is the absence of \texttt{sorryAx},
robust against the word \texttt{sorry} appearing inside \texttt{sorry-free} annotations
in the sources).

\subsection{Discharging the equality classification (D1--D4)}\label{ssec:lean-brouwer}
Theorem~\ref{thm:brouwer} has two parts: the edge \emph{bound} (Brouwer's
$e\le t_r(n)-\lfloor n/r\rfloor+1$ for non-$r$-partite $K_{r+1}$-free graphs) and the
\emph{equality classification} of the extremal graphs (Kang--Pikhurko). Both are now
formalised in Lean. The bound is \texttt{kp\_saving}: the full Kang--Pikhurko induction
--- Zykov-type symmetrisation, Case~A, the good/bad dichotomy, Lemma~3's
$K_{r+1}$-counting, and the singleton guard --- axiom-clean. The classification is
\texttt{kp\_equality\_classification\_proven}, the object this subsection documents.

The statement discharged is exactly the one Corollary~\ref{cor:equality21} consumes:
every graph $F$ on $21$ vertices with $\alpha(F)\le 5$, $K_5$-free and $e(F)=37$ is
isomorphic to a Kang--Pikhurko extremal construction --- equivalently its complement
$J=\overline F$, which is $K_6$-free, non-$5$-partite and has $173$ edges, is a maximum
such graph. This is the classification of Kang and Pikhurko~\cite{KP05} at the single
leaf $(r,n)=(5,21)$; we take that as the theorem \emph{statement} and reprove it in Lean
by a self-contained cone descent, in four steps.
\begin{itemize}[nosep]
\item[\textbf{D1}] \emph{Cone extraction.} For an extremal $J$, a maximum-degree vertex
$x$ has an independent complement $C=V\setminus N[x]$ joined completely to $N(x)$: the
Zykov symmetrisation of $J$ at $x$ is again $K_{r+1}$-free with $\alpha\le 4$, so
\texttt{kp\_saving} forces its edge count to equal $e(J)$, and pointwise degree equality
then makes $J$ a cone over $J[N(x)]$ (\texttt{cone\_extremal\_gen}).
\item[\textbf{D2}] \emph{Forced descent.} The same bound, applied at $r-1$ to
$J[N(x)]$, forces the removed block $C$ to have size exactly $4$ at each level (of the
feasible sizes only $c=4$ survives, margins $1$--$2$), peeling one independent $4$-set
per level: $(5,21)\to(4,17)\to(3,13)\to(2,9)$
(\texttt{descent\_21\_to\_17}, \dots, \texttt{descent\_13\_to\_9}).
\item[\textbf{D3}] \emph{Join-transport reassembly.} Each peeled level is recorded as a
cone \texttt{coneExtend}, which is functorial for graph isomorphism
(\texttt{coneExtend\_congr}); so an extremal $(5,21)$ graph is three nested cones over
its $(2,9)$ base --- i.e.\ $J=K_{4,4,4}\ast(\text{base})$ --- and an isomorphism of bases
lifts to an isomorphism of the reassembled graphs (\texttt{extremal21\_giso}).
\item[\textbf{D4}] \emph{Base classification.} Every $9$-vertex triangle-free graph with
$\alpha\le 4$, $e=17$ and $\Delta\le 4$ is isomorphic to one of exactly two graphs
(\texttt{base\_classification}). This is proved by pure local counting --- the degree
sequence is forced to $[2,4^8]$ or $[3,3,4^7]$; a degree-$2$ apex fixes the first, and
an independence-squeeze on the two degree-$3$ vertices fixes the second --- with
explicit \texttt{Equiv} witnesses, \emph{not} by graph enumeration.
\end{itemize}
Reassembling the two bases through the three cone layers yields exactly two isomorphism
classes for $J$, the constructions \texttt{kpG} (internal parameter $\abs{A^\ast}=2$)
and \texttt{kpG1} ($\abs{A^\ast}=1$; the naive $\abs{A^\ast}=3$ variant is isomorphic to
it). The transport consuming the classification (\texttt{equality21\_transport}) is
variant-agnostic, so the uniform $A/B$ structure of Corollary~\ref{cor:equality21}
follows for either class. As an independent cross-check, canonical-form (nauty)
enumeration confirms exactly these two classes at $(5,21)$ and the classification at
five nearby parameter families.

At the initial formalisation (repository tag \texttt{v1.0}) the result was
\emph{conditional} on this classification, carried as the hypothesis
\texttt{KPEqualityClassification}; the D1--D4 discharge above closed it the same day,
removing the last mathematical hypothesis. The reservation this leaves is not
mathematical but one of provenance: the discharge is a fresh, load-bearing
formalisation produced by the same AI pipeline as the rest of the development
(Section~\ref{sec:provenance}), and its D4 base step and construction witnesses rest on
\texttt{native\_decide} (the axiom profile below). A reader who would rather not trust
the new formalisation can audit the published classification of~\cite{KP05} directly and
consume it through the still-exported conditional theorem.

\subsection{Axiom profile}
Running \texttt{\#print axioms Erdos617.erdos\_617\_r5\_unconditional} reports exactly:
\[
  \texttt{propext},\quad \texttt{Classical.choice},\quad \texttt{Quot.sound}
\]
(the three standard Mathlib axioms), plus fourteen named \texttt{native\_decide}
reflection axioms --- \emph{seventeen} axioms in all. The fourteen are instances of
\texttt{Lean.ofReduceBool}: four enter through \texttt{primFacts} (the SAT certificates,
as before), and ten are the finite Kang--Pikhurko construction facts used by the D4
discharge --- two cone-isomorphism witnesses (\texttt{kpG\_giso\_cone3},
\texttt{kpG1\_giso\_cone3}) and the two $A/B$ complement structures (four axioms each).
Each is a decidable check on a fixed graph on at most $21$ vertices. \texttt{ofReduceBool}
is the same trust base as Lean's \texttt{bv\_decide} tactic: the kernel trusts the
compiled Boolean evaluation of a decision procedure on fixed input --- for the four SAT
facts, the LRAT-certificate checker; for the ten new facts, a direct adjacency
computation on the explicit constructions. No \texttt{sorryAx} appears, and no
mathematical hypothesis remains in the theorem's type. Thus the formal guarantee is
precisely: \emph{granting the \texttt{native\_decide} evaluation of those fourteen fixed
checks, Theorem~\ref{thm:main} holds, with everything else --- including Brouwer's bound
and the equality classification --- checked by Lean's kernel.} This is a machine-checked
proof modulo that reflection trust base --- unconditional in that it assumes no
unformalised mathematical input, not a claim of a kernel-only proof. The still-exported
conditional \texttt{erdos\_617\_r5} carries only the three standard axioms and the four
SAT reflections, with \texttt{KPEqualityClassification} visible in its type.

\subsection{How to check}
Build with \texttt{lake exe cache get \&\& lake build} in \texttt{lean617/}; inspect
the guarantees with \texttt{\#print axioms Erdos617.erdos\_617\_r5\_\allowbreak
unconditional} (which should show the three standard axioms, the fourteen
\texttt{native\_decide} reflections, and no \texttt{sorryAx}, with no mathematical
hypothesis in the theorem's type). Comparing with \texttt{\#print axioms
Erdos617.erdos\_617\_r5} isolates the ten construction axioms the discharge adds. The
DRAT path~(i) of Section~\ref{sec:certificates} is a solver-independent cross-check of
the four SAT facts that does not rely on \texttt{native\_decide}.

%% ===========================================================================
\section{Provenance, verification, and limitations}\label{sec:provenance}
%% ===========================================================================

This work was produced with heavy use of AI systems. Following the norms of
\cite{Bloom617}, we disclose their role in full and give a candid account of what
has and has not been verified.

\subsection{Status of this report}
We state the framing plainly. The mathematics in this document --- both lemma proofs,
the reduction, the toolkit, the certificates, and the Lean formalisation --- was
\emph{produced and internally reviewed by AI systems}. The human owner of this effort
has verified the \emph{process} (that the reduction is sound, that the reviews were
independent and adversarial, that the certificates check, that the Lean build is
clean with the stated axiom profile) but does \emph{not} assert having personally
verified the mathematical content of the two lemma proofs, and is not posting a
personal claim that the theorem is established. The document is therefore offered as
a verification-ready report for independent mathematician review, not as an authored
announcement. Section~\ref{sec:howto} is the guide to performing that review.

\subsection{Who did what}
\begin{itemize}[nosep]
\item The proofs of Lemma~\ref{lem:MH} (\textsf{[MH$''$]}) and Lemma~\ref{lem:MM}
(\textsf{[MM]}) --- the two hard combinatorial arguments of
Sections~\ref{sec:MH}--\ref{sec:MM} --- were drafted by OpenAI's
\texttt{gpt-5.6-sol} (via the \texttt{codex} CLI) from self-contained briefs, in the
form presented in the repository's \texttt{review\_queue/}.
\item Everything else was produced by Claude (Anthropic) sessions: the problem
framing and the reduction (Lemma~\ref{lem:chain}); the empirical structure-mining
that identified the correct lemma statements (Appendix~\ref{app:phase}); the briefs
handed to \texttt{gpt-5.6-sol}; the \capeleven{} toolkit exposition; the SAT
encodings, certificates, and their re-verification; and the entire Lean formalisation
--- the reduction, both lemmas, Brouwer's bound (\texttt{kp\_saving}), and the D1--D4
discharge of the equality classification (Section~\ref{ssec:lean-brouwer}), this last
carried out by a relay of Claude sessions --- together with the adversarial reviews
described next. This write-up was assembled by a Claude session from those frozen
sources.
\end{itemize}

\subsection{The adversarial-review chain}
Each of the three links was reviewed by a \emph{fresh} model session instructed to
attack it, working independently of the author session and re-deriving or
re-checking every load-bearing step.
\begin{itemize}[nosep]
\item The reduction (Lemma~\ref{lem:chain}) was reviewed and \textsc{accepted}: all
six steps and both lemma interfaces were found sound.
\item Lemma~\ref{lem:MH} was \textsc{accepted modulo} one repairable gap: the
original draft justified a clique bound $\omega(H[L])\le 4$ that the reviewer could
not reconstruct. The repair --- used here in Section~\ref{ssec:MH-K5free} --- replaces
it with the \emph{$\omega$-free} $11$-vertex nonexistence (Fact~\ref{sat:nonex11}),
which makes the offending case vacuous outright. The reviewer independently
re-derived the counting identities, every $M$- and $L$-value, the affine bounds and
slacks, and the finite endgame facts.
\item Lemma~\ref{lem:MM} was \textsc{accepted modulo} one missing conclusion: the
$\nu=7$ case of Section~\ref{ssec:MM-two} stopped at the forced weight vector
$(0,0,0,6,6)$ without closing. The reviewer supplied the closing paragraph (reusing
the $\nu=4$ argument), which we adopt here. The reviewer re-ran the $g$-table, the
case scan~\eqref{eq:MM-16}, and the finite graph facts on independent encodings.
\item Theorem~\ref{thm:affine} was reviewed and \textsc{accepted}, with an
independent exhaustive check at $q=3$.
\end{itemize}
The repaired proofs are the ones presented in this paper. In addition, the one
external theorem (Theorem~\ref{thm:brouwer}) was checked against the primary sources
\cite{KP05} and a secondary statement of Brouwer's bound~\cite{RWWY}, with all five of its uses
confirmed faithful and all Tur\'an arithmetic recomputed.

\subsection{What has and has not been verified}
\emph{Verified (by us):} the reduction and both lemmas have each been (a) drafted or
constructed, (b) checked line-by-line by the author session, (c) subjected to an
independent adversarial review that re-derived the load-bearing content, and
(d) formalised in Lean with no \texttt{sorry}s and no remaining mathematical
hypothesis --- Brouwer's bound and the Kang--Pikhurko equality classification at
$(5,21)$ included (Section~\ref{sec:lean}). Every finite numerical claim is confirmed
by exact
recomputation; every SAT fact carries a DRAT certificate checked by
\texttt{drat-trim} \emph{and} an LRAT certificate checked inside Lean; the external
theorem is literature-verified.

\emph{Not yet verified:} the chain of author sessions, fresh-session adversarial
reviewers, and machine checks is an \emph{internal} verification. It has not been
refereed by an independent human expert, which for a result of this kind is the
appropriate next standard. Two features of the formal object a verifier should weigh
remain: (i) the discharge of the equality classification
(Section~\ref{ssec:lean-brouwer}), while it removes the last mathematical hypothesis,
is itself a fresh, load-bearing formalisation by the same AI pipeline, reviewed only
internally --- a reader who distrusts it can instead audit the published classical
statement and use the still-exported conditional theorem; and (ii) the SAT facts and
the finite Kang--Pikhurko construction facts are trusted through
\texttt{native\_decide}'s reflection axiom (with the independent \texttt{drat-trim} path
as mitigation for the SAT facts). We make no claim that the present internal
certificate substitutes for expert human scrutiny; we invite it.

\subsection{Why the argument should be believed, and where to attack it}
The proof is elementary and self-contained modulo one classical theorem, and it uses
the hypotheses essentially: it distinguishes $r^2$ from $r^2+1$ (it fails, as it
must, at $n=24$; Section~\ref{ssec:MH-n24} and Appendix~\ref{app:phase}), it does not
prove too much (the pure-graph relaxations of both lemmas are false), and it is not
suspiciously short. The most fragile-looking points --- and the ones a sceptical
reader should probe first --- are the $n=21$ equality classification feeding
Corollary~\ref{cor:equality21} (now formalised by the D1--D4 cone descent of
Section~\ref{ssec:lean-brouwer}, whose $(2,9)$ base step and \texttt{native\_decide}
construction witnesses are the natural targets), the correctness of the $L$-table
recursion (whose affine bounds we verify pointwise), and the two repaired cases
(Sections~\ref{ssec:MH-K5free} and~\ref{ssec:MM-two}) that turn on the $\omega$-free
$11$-vertex nonexistence. Each of these is exactly where an independent check would
be most valuable.

%% ===========================================================================
\section*{Acknowledgements}
%% ===========================================================================
\addcontentsline{toc}{section}{Acknowledgements}
We thank the maintainers of \cite{Bloom617} and of the \texttt{formal-conjectures}
project~\cite{FC} for the infrastructure that makes a claim like this one checkable,
and the authors of kissat, cadical, \texttt{drat-trim}, Lean, and Mathlib.

%% ===========================================================================
\appendix
\section{The empirical phase transition}\label{app:phase}
%% ===========================================================================

The lemma statements were not guessed; they were extracted from computation, which
also explains why the argument bites at $n=25$ and not below. We record the data
here as motivation and as a falsification guard.

\emph{Balanced $K_{25}$ colourings are plentiful and rigid.} Local search finds
balanced $5$-colourings of $K_{25}$ routinely from random starts; over $120$
distinct referee-verified samples were generated, including non-affine ones (edge
profile $(52,54,58,68,68)$ versus the affine $(50,50,50,50,100)$). Uniformly across
all samples: every class has exactly $5$ monochromatic $K_5$'s and hitting number
$5$, and every minimal hitting $5$-set spans $9$ or $10$ edges of its own class ---
never $\le 6$, the ``usability'' threshold of Step~3 of Lemma~\ref{lem:chain}. Every
sample fails one-vertex extension (all $\binom{25}{6}$-set checks pass for $K_{25}$
but the extension CNF is unsatisfiable).

\emph{The hitter family dies between $24$ and $25$.} Call an \emph{h4-witness at $n$}
a balanced $5$-colouring of $K_n$ with a colour and a $4$-set $T$ such that
$\alpha(G_0-T)\le 4$ --- exactly the configuration Lemma~\ref{lem:MH} forbids at
$n=25$. Such witnesses \emph{exist} and are referee-verified for every
$n\in\{17,\dots,24\}$ (found by SAT for $n\le 22$ and by local search for $n=23,24$).
At $n=25$, five independent searches each descend to exactly $24$ residual violations
and stall there for tens of CPU-hours, with every residual violation passing through
the newest vertex. This ``deficiency-$1$'' signature --- the hitter constraint costs
exactly one vertex of slack --- is the empirical shadow of the ``$+1$'' in the
conjecture, and matches the structural fact that $21=4\cdot 5+1$ forces one oversized
part in the proof of Lemma~\ref{lem:MH} (Section~\ref{ssec:MH-n24}) whereas
$20=4\cdot 5$ splits evenly.

\emph{Pure-graph relaxations fail.} The graph relaxation of Lemma~\ref{lem:MH}
(drop balancedness, keep $\alpha\le 4$ and \capeleven{} on $21$ vertices) has an
$80$-edge witness; the graph relaxation of Lemma~\ref{lem:MM}
($\alpha\le 5$, \capeleven, $\le 60$ edges on $25$ vertices) is satisfiable. Both
lemmas therefore genuinely require the surrounding colouring structure that the
reduction supplies --- which is why the proof cannot be replaced by a single Tur\'an
computation, and why the machine facts it does use are confined to graphs on at most
$16$ vertices.

%% ===========================================================================
\begin{thebibliography}{9}
%% ===========================================================================

\bibitem{Bloom617}
T.~F. Bloom,
\emph{Erd\H{o}s Problem \#617},
\url{https://www.erdosproblems.com/617}, accessed 2026-07-11.

\bibitem{Brouwer81}
A.~E. Brouwer,
\emph{Some lotto numbers from an extension of Tur\'an's theorem},
Math. Centrum report ZW~152, Mathematisch Centrum, Amsterdam, 1981.

\bibitem{ErGy99}
P.~Erd\H{o}s and A.~Gy\'arf\'as,
\emph{Split and balanced colorings of complete graphs},
Discrete Math. \textbf{200} (1999), 79--86.
\url{https://doi.org/10.1016/S0012-365X(98)00323-9}.

\bibitem{FC}
Google DeepMind,
\emph{formal-conjectures} (Lean formalisations of open problems),
\url{https://github.com/google-deepmind/formal-conjectures}
(file \texttt{FormalConjectures/ErdosProblems/617.lean}), accessed 2026-07-11.

\bibitem{KP05}
M.~Kang and O.~Pikhurko,
\emph{Maximum $K_{r+1}$-free graphs which are not $r$-partite},
Mat. Studii \textbf{24} (2005), no.~1, 12--20.
\url{https://doi.org/10.30970/ms.24.1.12-20}.

\bibitem{RWWY}
S.~Ren, J.~Wang, S.~Wang, and W.~Yang,
\emph{Extremal triangle-free graphs with chromatic number at least four},
preprint, arXiv:2404.07486 (2024).
\url{https://arxiv.org/abs/2404.07486}.

\end{thebibliography}

\end{document}
